\documentclass[11pt,a4paper]{article}
\usepackage[utf8]{inputenc}
\usepackage[T1]{fontenc}
\usepackage{lmodern,amsmath,amssymb,amsthm,mathtools,booktabs,enumitem,microtype}
\usepackage{natbib}
\usepackage[margin=1in]{geometry}
\usepackage{xr-hyper}
\usepackage{hyperref}
\externaldocument[PC-]{01i_graphic_equalisation_passive_continuation_momentum_gravity}
\externaldocument[HG-]{01h_graphic_equalisation_information_gravity_quantum_reconciliation}
\externaldocument[GE-]{01_law_of_graphic_equalisation_theory_referenced}
\externaldocument[IS-]{00_is}
\externaldocument[MS-]{01b_graphic_equalisation_mind_and_self_grounding}
\externaldocument[NM-]{03_law_of_narrative_momentum_theory_referenced}
\externaldocument[BN-]{01d_graphic_equalisation_born_recovery}
\newtheorem{definition}{Definition}[section]
\newtheorem{theorem}[definition]{Theorem}
\newtheorem{proposition}[definition]{Proposition}
\newtheorem{corollary}[definition]{Corollary}
\theoremstyle{remark}\newtheorem{remark}[definition]{Remark}
\title{The Law of Graphic Equalisation --- Physics Grounding\\\large Change, $c$, the Reconciliation of Special and General Relativity, the Recovery of the Born Rule, the Dimensional Constants, and the Second Law}
\author{Matthew Parslow\\Independent Researcher}
\date{\textbf{Draft \texorpdfstring{\textemdash}{---}\ \PaperDraft}\\[2pt]\PaperDisplayDate \texorpdfstring{\textbullet}{\textperiodcentered}\ document revision \PaperRevision}
\usepackage{fancyhdr}
\pagestyle{fancy}
\fancyhf{}
\fancyfoot[C]{\thepage}
\fancyfoot[R]{\scriptsize\textsc{\PaperDraft\,r\PaperRevision}}
\renewcommand{\headrulewidth}{0pt}
\fancypagestyle{plain}{\fancyhf{}\fancyfoot[C]{\thepage}%
  \fancyfoot[R]{\scriptsize\textsc{\PaperDraft\,r\PaperRevision}}\renewcommand{\headrulewidth}{0pt}}
% GENERATED by tools/build_paper_metadata.py from paper_metadata.json.
% DO NOT EDIT: edit paper_metadata.json, then regenerate this file.
\providecommand{\PaperDraft}{draft-2026-09-06.1}
\providecommand{\PaperDisplayDate}{6 September 2026}
\expandafter\def\csname PaperRevision@00_is\endcsname{129}
\expandafter\def\csname PaperRevision@01_law_of_graphic_equalisation_theory_referenced\endcsname{294}
\expandafter\def\csname PaperRevision@01a_graphic_equalisation_physics_grounding\endcsname{47}
\expandafter\def\csname PaperRevision@01b_graphic_equalisation_mind_and_self_grounding\endcsname{47}
\expandafter\def\csname PaperRevision@01c_graphic_equalisation_cosmology_grounding\endcsname{11}
\expandafter\def\csname PaperRevision@01d_graphic_equalisation_born_recovery\endcsname{6}
\expandafter\def\csname PaperRevision@01e_graphic_equalisation_communication_grounding\endcsname{6}
\expandafter\def\csname PaperRevision@01f_graphic_equalisation_qualia_and_consciousness_grounding\endcsname{3}
\expandafter\def\csname PaperRevision@01g_graphic_equalisation_total_perspective_vortex_grounding\endcsname{4}
\expandafter\def\csname PaperRevision@01h_graphic_equalisation_information_gravity_quantum_reconciliation\endcsname{5}
\expandafter\def\csname PaperRevision@01i_graphic_equalisation_passive_continuation_momentum_gravity\endcsname{12}
\expandafter\def\csname PaperRevision@01s_graphic_equalisation_statement\endcsname{294}
\expandafter\def\csname PaperRevision@02_seldon_manifolds_theory_referenced\endcsname{31}
\expandafter\def\csname PaperRevision@02a_seldon_manifolds_psychological_grounding\endcsname{15}
\expandafter\def\csname PaperRevision@02b_seldon_manifolds_social_and_psychological_application\endcsname{7}
\expandafter\def\csname PaperRevision@02c_seldon_manifolds_decision_making\endcsname{11}
\expandafter\def\csname PaperRevision@03_law_of_narrative_momentum_theory_referenced\endcsname{32}
\expandafter\def\csname PaperRevision@04_narrativium_theory_referenced\endcsname{23}
\expandafter\def\csname PaperRevision@05_frames_relation_and_the_empty_limit\endcsname{11}
\expandafter\def\csname PaperRevision@06_assessing_an_account_of_this_kind\endcsname{29}
\expandafter\def\csname PaperRevision@07_causal_attitude_networks_on_the_substrate\endcsname{4}
\expandafter\def\csname PaperRevision@08_certified_answers_as_an_audit\endcsname{4}
\expandafter\def\csname PaperRevision@09_adaptability\endcsname{94}
\expandafter\def\csname PaperRevision@09a_adaptability_implementation\endcsname{7}
\expandafter\def\csname PaperRevision@10_applied_graphic_equalisation\endcsname{7}
\expandafter\def\csname PaperRevision@10a_red_queen_race_civilisational_rate_of_change_llms_biological_timescales_anthropocene_transition\endcsname{13}
\expandafter\def\csname PaperRevision@11_relational_information_ethics\endcsname{16}
\expandafter\def\csname PaperRevision@A1_constructive_graphic_equalisation_implementation_referenced\endcsname{8}
\expandafter\def\csname PaperGenerated@01s_graphic_equalisation_statement\endcsname{1}
\edef\PaperMetadataStem{\jobname}
\expandafter\ifx\csname PaperRevision@\PaperMetadataStem\endcsname\relax
  \PackageError{paper_metadata}{No metadata for \PaperMetadataStem}{Add the paper stem to paper_metadata.json.}
  \def\PaperRevision{?}
\else
  \expandafter\let\expandafter\PaperRevision\csname PaperRevision@\PaperMetadataStem\endcsname
\fi
\newcommand{\PaperDateLine}{\textbf{Draft \texorpdfstring{\textemdash}{---}\ \PaperDraft}\\[2pt]\PaperDisplayDate\ \texorpdfstring{\textbullet}{\textperiodcentered}\ document revision \PaperRevision}
\newcommand{\PaperStatementDateLine}{\textbf{Draft \texorpdfstring{\textemdash}{---}\ \PaperDraft}\\[2pt]\PaperDisplayDate\ \texorpdfstring{\textbullet}{\textperiodcentered}\ statement}
\expandafter\ifx\csname PaperGenerated@\PaperMetadataStem\endcsname\relax
  \newcommand{\PaperFooterStamp}{\textsc{\PaperDraft\,r\PaperRevision}}
\else
  \newcommand{\PaperFooterStamp}{\textsc{\PaperDraft\ / statement / from long r\PaperRevision}}
\fi
\newcommand{\PaperNoticeStamp}{\textbf{DRAFT --- \PaperDraft\ r\PaperRevision} (\PaperDisplayDate).}


\begin{document}
\maketitle
\begin{center}
\fbox{\parbox{0.92\textwidth}{\small
\textbf{DRAFT --- \PaperDraft r\PaperRevision} (\PaperDisplayDate).
This document is under active revision and is circulated for comment, not as a
finished result. It is a companion to \emph{The Law of Graphic Equalisation ---
Theory} (referred to below as Graphic Equalisation), and inherits that paper's
definitions, epistemic boundary, and level-of-claims. The other documents of the
programme named below carry short names too, and those are the only forms used:
\emph{Is}, the mind companion (\emph{\ldots\ Mind and Self
Grounding}), the momentum paper, and the Born note (\emph{\ldots\ Born
Recovery}). Sections marked as
falsifiability commitments are the intended points of attack.}}
\end{center}
\vspace{1em}

\begin{abstract}
This companion explores conditional physics interpretations of Graphic Equalisation.
A recurrent frame can normalise its clock rate to one; relating that convention to
observer-independent light propagation requires additional physical assumptions.
Relativistic proper time depends on the spacetime metric and the clock's worldline;
a fixed change-budget is a proposed interpretation, not a derivation of that metric
or a reconciliation of special and general relativity. Quantum amplitudes, measurement,
and entanglement are compared with frame-relative structures. Recovering the Born
rule requires a Hilbert-space or operational-probabilistic grounding whose additivity,
noncontextuality, continuity, purification, and composition premises are not all
established here. Landauer's bound concerns logical erasure, and the second law
requires thermodynamic assumptions beyond overwriting. Indistinguishability does
not by itself imply fermionic exclusion. The paper derives no field equation,
dimensionless constant, or theory of quantum gravity; its proposed correspondences
are accompanied by assumptions and remaining proof obligations.
\end{abstract}

\section{Scope}\label{sec:scope}
This paper is a grounding, not a physical theory. It inherits Graphic
Equalisation's position that the claims are about information --- distinction,
frame, admissibility, realisation --- of which physical theory is one embodiment,
and that a physical result bears on an information-level claim when it settles an
informational fact. Everything below is at that level. The proposed correspondences are structural interpretations. A clock-unit convention
alone supplies neither Lorentz invariance nor gravitational redshift. Born-rule
uniqueness is conditional on the hypotheses of the cited reconstruction theorems.
The paper supplies no field equation, quantitative metric, or quantum-gravity model.
A failed correspondence can therefore identify an inadequate grounding without
refuting the abstract substrate definitions.

The minimal-universe argument that licenses the contingency verdicts below --- a
rock is a universe, and an inert coarse description can omit succession without thereby proving that physical gauge structure, constants, or spatial dimensions are absent --- together with the
run/record axis, free will, personal identity, and death, is developed in the
mind companion (its Section~\ref{MS-sec:minimal} and
Proposition~\ref{MS-prop:minimal-universe}), which treats philosophy of mind as
applied physics of the run.

\section{\texorpdfstring{$c$}{c} as the change-to-time rate}\label{sec:c-rate}
Graphic Equalisation establishes that time is not a substrate dimension but is
\emph{inferred} by a frame. This grounding provisionally uses a closure cycle --- a recurrence
modelled as one traversal of a closed rule-loop --- as its operational unit, of no fixed hidden
length. That is a simplest working clock model, not a proof that the substrate closes or supplies a
uniform cycle. A clock may be externally supplied, or a superframe may choose when a contained frame
is advanced, while every clock accessible inside the contained frame is rescheduled with it.

Conditional on this closure-driven realization, the invariant speed $c$ has a natural reading,
consistent with the missing-denominator result (Graphic Equalisation's Remark~\ref{GE-rem:no-denominator}: there is no
rate \emph{in} substrate-time, since there is no substrate-time to divide by). $c$ is then a
normalization of represented change-throughput against the frame's accepted clock unit, not a direct
reading of a substrate rate. That is \emph{runner} $c$, the first of the three things the word names
here, and the separation has to come before anything is claimed of any of them.

A recent collapse-model calculation reaches the empirical pressure point by a different route.
Within the Diósi--Penrose and continuous-spontaneous-localisation models, \citet{bortolotti2025}
take a stochastic Newtonian potential over an already supplied time parameter and derive a
corresponding uncertainty in a clock reading. That is a conditional, quantitative result about
those models --- not evidence that this substrate is established, and not a derivation of time:
the fluctuating potential, its correlation function, and the parameter with respect to which it
fluctuates are model inputs. Its relevance here is narrower and independent. It confirms that a
physical account of clocks can make their temporal rate and precision depend on physical
conditions rather than treating an exact background time as all a clock reads. The question there
is why a clock reading fluctuates \emph{given} time; the question here is what a clock is counting
when it represents time at all. In the present closure-driven grounding the latter is the frame's
accepted recurrence count, whose variable change-load supplies the proposed reason a clock rate can
vary; that realization is provisional rather than a claim that every substrate closes.

\begin{definition}[Three things called $c$]\label{def:three-c}
Three things are called $c$ in this paper and only the third is empirical, while
only the first is what the substrate is about. \emph{Runner $c$} names whatever rate or schedule supplies the frame; it is never
readable from inside, and no route described here reaches or proves its regularity. \emph{Unit $c$}
is $1$ by local normalization in the frame's accepted operational clock unit, closure-based by
default in this grounding; it is definitional and carries no empirical risk. \emph{Observed $c$} is
the number a frame obtains by measurement, retention and inference through whatever counters its rules
happen to provide. Only the normalization is fixed here; runner scheduling and observed values remain
contingent rule content.
\end{definition}

\begin{proposition}[Unit $c$ is each frame's local normalization]\label{prop:c-unit}
A frame normalizes its accepted operational clock unit to $1$ --- not because
anything has been measured, but because it holds no independent substrate duration
against which to read that unit. In the closure-driven realization this is a
closure-count; in a clock-driven realization it may be an externally supplied
counter. The invariance of \emph{unit} $c$ is therefore definitional: not that all
frames agree on a speed or scheduling rate, but that each reports its own unit.
Nothing is thereby claimed of runner $c$, which is not readable, or of observed
$c$, which is contingent (Definition~\ref{def:three-c}).
\end{proposition}

\begin{remark}[Provisional closure, external clocks, and superframe scheduling]\label{rem:clock-provisional}
A silicon computation illustrates the distinction. Its logical step may be bounded by a clock edge
supplied from outside the program; a dedicated substrate could instead advance locally when a rule
settles. Either is compatible with the same internal succession. More generally, a containing frame
may choose when to grant a contained frame another executable step, batch steps, pause it, or vary the
hidden work between them. If every counter available to the contained frame is rescheduled together,
the frame cannot decide from inside whether its apparent recurrence is self-closing, clocked, or
superframe-scheduled. Closure is therefore the shortest working explanation where no contrary evidence
is held, not a required substrate fact.
\end{remark}
A value \emph{other} than $1$ is a cross-frame quantity: the familiar number
appears only when one frame is compared with another --- change-throughput
against a spatial metric, seconds against metres --- which is where the
continuous lives, in the maps between frames rather than in a frame's own
discrete states. One can give time a number at all only by comparing frames; in
its own frame $c$ is the unit, and its particular value ($299\,792\,458$ in
SI) is an artefact of the comparison, not a fact about the substrate. And $c$ is
a \emph{limit} because a signal is a sequence of changes and none can be realised
faster than the frame's throughput.

Local normalisation fixes a reported unit, not the physical transformation
between frames. Recovering observed Lorentz invariance needs a shared signal
law and a transformation rule. The budget model below imports the relativistic
quadratic relation as a grounding assumption. It is not derived by counting
closure events, and a local cycle convention alone predicts neither a universal
signal speed nor gravitational clock rates.

\section{The change-budget: speed, dilation, mass, aperture}\label{sec:budget}
For the conditional closure-driven realization developed below, the budget is modelled as
throughput per accepted clock unit, and $c$ is its \emph{unit} rather than its size
(Definition~\ref{def:three-c}). This is a grounding assumption, not a claim that every substrate
has a fixed intrinsic throughput per self-closing cycle. That fixed throughput splits between
\emph{spatial motion} and \emph{internal change} (proper time) --- in squares, the
two shares summing to one, which is the constant-magnitude four-velocity below ---
and a speed is the share spent on motion. At rest the whole budget is internal
change and the clock ticks fully; increasing speed moves budget to motion, leaving
less for internal change, so the clock slows --- time dilation as budget
re-allocation --- and at $c$ the whole budget is motion and no proper time remains.
Mass is what holds the internal share \emph{above zero at every finite speed}: a
massive graph's motion share approaches one and never reaches it, so it asymptotes
to $c$ and the energy to approach diverges, while a massless one has no internal
share at all and \emph{is} at $c$, with no proper time.

Because observation is funded by internal-change throughput, approaching $c$
coarsens a frame's \emph{aperture}. The transformation between frames in relative
motion --- time dilation and length contraction --- is an aperture transformation
between different budget-splits, and each frame's aperture self-calibrates so that
its own $c$ still reads $1$. The Lorentz factor is the aperture change that holds
$c=1$ across the split. (The constant-magnitude four-velocity of Minkowski
geometry \citep{minkowski1908} is this budget-split stated as a metric fact
rather than as a mechanism.)

\begin{remark}[$c$ is not directly readable from inside, and knowing it costs
retention]\label{rem:c-indirect}
Being $1$ by construction has an epistemic consequence that is easy to miss and belongs with the
rest. Well-formedness does not guarantee opaque provenance (Graphic Equalisation's Remark~\ref{GE-rem:exposure-leakage}). Here $c$ is the rate at
which whatever supplies them supplies them, as the frame meets it. So $c$ is not a datum a frame can read off directly: any in-frame attempt to measure it is itself
funded by the same throughput, and returns $1$ whatever the underlying rate --- a clock cannot time
its own tick. That is exactly why the unit-reading carries no risk and no content.
\end{remark}

\begin{remark}[The indirect route is available, and it is priced]\label{rem:c-inferred}
It does not follow that $c$ is unknowable, and the opposite end has to be stated with equal
firmness, because it has been exercised for three centuries. A frame with measurement, retention and
inference obtains $c$ \emph{indirectly}, as a number in its own other units --- Rømer from eclipse
timings held across months, Fizeau's toothed wheel and the rotating-mirror
measurements after it from distance set against duration ---
and that route is available to any frame that can hold a predecessor and compare. The two ends are
not in tension: what is absent is the direct read, and what is present is the inference, exactly as
with a frame's own aperture limit, which leaves no direct residual and is nevertheless inferable ---
weakly, from imposed change the frame cannot attribute --- as the mind companion's
Proposition~\ref{MS-prop:selfmodel} sets out.

Two things follow. Knowing $c$ has a price, and it is the one the write-economics already names: the
inference needs at least two events and a retained record of the first, and retention is the paid
copy (Graphic Equalisation's write-economics, Proposition~\ref{GE-prop:economics}). A frame with no
memory does not have a
poor estimate of $c$; it has no access to the question. And the indirection is what makes the claim
\emph{empirical} at all. If $c$ were directly readable it would be definitional and unfalsifiable;
because the only route to a number runs through measurement and inference, a disagreement could in
principle show up there, which is where the falsifiable content of Section~\ref{sec:falsifiability}
sits --- not in $c=1$, but in the claim that two frames taking that indirect route cannot be
exhibited reaching different answers for one signal.

This is also, notably, not a proposal. Metrology reached the same conclusion and acted on it: since
1983 the metre has been \emph{defined} from $c$ rather than $c$ measured in metres, and the 2019
revision kept $c$ among the defining constants \citep{bipm2019}. A quantity whose value is fixed by
definition is not a fact about the world that experiment reports; it is a conversion between a
frame's length and time units. What the substrate adds is not that convention but a reason it could
be adopted without loss.

What this remark has been about throughout is the third of the three
(Definition~\ref{def:three-c}): runner $c$ is unreachable and unit $c$ is fixed, so everything
empirical sits in observed $c$. That is contingent rule content, and the following remark is about
what the contingency permits.
\end{remark}

\begin{remark}[The one-way speed of light is not a hard measurement but an absent
quantity]\label{rem:oneway}
The immediate consequence is a known result, and it is worth being clear at the outset that the
result is not this paper's. The one-way speed of light cannot be determined without first stipulating
a synchronisation for spatially separated clocks, and any such stipulation already fixes the answer;
what is measurable, and measured to great precision, is the two-way speed over a closed path. The
freedom in the choice is Reichenbach's $\varepsilon$, and the equivalence of the resulting theories
has been worked out in detail \citep{reichenbach1958,anderson1998}. Nothing below revises that. What
the substrate supplies is a reason for it that is not specific to light.

The reason, on this grounding's closure-clock hypothesis, is the chosen operational unit. A round
trip closes relative to one retained record and needs no agreement with anything elsewhere; a one-way
path does not. The result is not a theorem that no other counter can exist: an external clock,
unlocked in-frame counter, or superframe relation can alter what a frame can operationally compare.
Rather, where the frame has only its closure-based unit, no one-way duration is derived from that unit.
The obstruction is therefore conditional on the represented clock structure, not a substrate claim.

It is the same obstruction as the one above, arrived at from the other side. To time an open path a
frame needs two separated clocks that agree on simultaneity, and establishing that agreement requires
sending something between them at the rate being measured. The apparatus presupposes its own result,
exactly as an in-frame measurement of $c$ returns $1$ because it is funded by the throughput it is
trying to read.
\end{remark}

\begin{remark}[The inaccessibility is rule content, not a theorem]\label{rem:oneway-rule}
The qualification that makes this contingent rather than necessary has to be stated, because the
argument as given proves more than it is entitled to. Nothing forbids rules under which some
dimension steps $N$ times per closure cycle, and those steps would be perfectly observable from
inside --- they are in-frame events, not hidden variables. A frame holding such a dimension has a
counter finer than its own cycle, and if that counter also runs at a common rate across separated
locations it \emph{is} a global simultaneity, in which case the one-way speed is measurable in that
universe and the argument above does not apply to it. So the inaccessibility is a fact about rule
content, not a theorem about frames. What can be said generally is narrower and survives: a frame
cannot establish that any counter it holds is faithful to the rate beneath its primitives. Provenance
is not sealed --- primitives leak, and inference from in-frame observables is exactly how a number for
$c$ was ever obtained (Graphic Equalisation's Proposition~\ref{GE-prop:leak}) --- but every such
inference is
itself performed by primitives whose provenance is open in the same way, so the check relocates
rather than closing. Where the counter is locked to the closure cycle, observed $c$ tracks unit $c$
and nothing is measurable one-way; where it is not, observed $c$ may differ from the cycle rate, by
location or by direction, and the frame has no way to tell which case it is in. \emph{Observed $c$ is
always a report about a frame's counters, and reaches what supplies its primitives only by inference
from them and never by reading it}. That much holds without qualification; what does not is the
stronger claim that such inference is unavailable --- it is available,
it is what three centuries of measurement did, and it constrains without settling.
\end{remark}

\begin{remark}[The locked case is the cheaper hypothesis, which is not a proof]\label{rem:locked-cheaper}
There is an economic reason to expect the locked case without needing to legislate it. A dimension
stepping at a common rate everywhere is structure the substrate must maintain across every location,
and maintenance is paid (Graphic Equalisation's Proposition~\ref{GE-prop:economics}); the closure cycle already supplies a
unit at no such cost. The locked case is in that sense the \emph{better-formed} one --- fewer parts,
nothing held that does no work, the same well-formedness Graphic Equalisation's cost result turns on
(Proposition~\ref{GE-prop:economics}).
It is therefore both the cheaper hypothesis and the shorter description, and a frame with no direct
access should assume it on those grounds \citep{rissanen1978}. Such dimensions are accordingly
expected to be rare, which is not a reason they cannot occur; our own universe's apparent lack of one
is an empirical finding rather than a derived necessity.

That licenses the assumption without forcing the fact, and the distinction is the whole of it. Being
better formed makes something likely, expected, and rational to assume; it does not make it
necessary, and there is no step from the one to the other. It is also why the standard
synchronisation feels non-arbitrary rather than merely stipulated --- it is the minimum-description
choice, which is exactly the kind of thing that gets adopted universally without ever being proved.
The failure to guard against is therefore not the assuming, which is correct practice. It is the
promotion of a well-founded assumption to a theorem.
\end{remark}

Both ends need stating, because the negative half is the smaller half. A frame learns a great deal
here: the two-way speed, its isotropy, and the constraints that isotropy places on any anisotropic
account are all in hand and tightly bounded by experiment. The sharpest item on that side is one-way
after all, and it is worth spelling out because it reads at first like a counterexample. The
\emph{absolute} one-way speed is unavailable; the \emph{ratio} of one-way propagation times through two media is
not, and is measured as a matter of routine --- a pulse sent along a fixed path through a medium
arrives late against the same path in vacuum by an amount that is a real datum. That media slow light
by large factors is itself routine: the group velocity in an ultracold gas has been reduced to
seventeen metres per second and read off \citep{hau1999}. The synchrony convention cancels: whatever
offset $\varepsilon$ assigns to the far end is a function of the separation, both signals cross the
same separation, so the offset is common to them and absent from the difference
\citep{anderson1998}. The condition carrying that is \emph{along the same path} --- two one-way
measurements in different directions do not share an offset under an anisotropic convention, and
their ratio is convention-dependent again.

\begin{proposition}[Closure sorts three cases]\label{prop:closure-three}
The closure criterion sorts three cases rather than two, and sorts them by one rule. The two-way
speed closes and is measurable; the absolute one-way speed does not close and is not; the one-way
ratio closes at the comparison though neither leg does, and is. A ratio is dimensionless and so never
needed expressing in the frame's cycle at all, and the comparison is local --- two arrivals set
against one record at one place, which closes.
\end{proposition}

This is the closure criterion delivering the positive result rather than an exception to it, and
getting all three cases from the same criterion is a check the account passes rather than a
convention it adopts. And the impossibility is indexed to frames \emph{inside} the structure being
measured, which is Axiom~$-1$ in its plainest form. That qualification has to be read as a real one
and not as decoration. The framework supplies outsides routinely: a superframe is one, and
Section~\ref{sec:entanglement}'s reading of entanglement as a superframe dimension is load-bearing
for the recovery of the Born rule two sections later. \emph{Not measurable in frame} is therefore the
whole of the claim, and it does not extend to \emph{not measurable}.

\begin{remark}[Reichenbach and Malament are the two halves of one situation]\label{rem:reichenbach-malament}
The strongest reply is worth engaging rather than omitting, and it is congenial. Standard synchrony
is not merely one convention among equals: it is the only nontrivial simultaneity relation definable
from causal connectibility together with an inertial worldline \citep{malament1977}, so
$\varepsilon=\tfrac12$ is canonical rather than arbitrary. That does not restore measurability of the
one-way speed, and it is not offered against the reading here --- causal connectibility \emph{is} the
round trip, so Malament's derivation and the cycle unit privilege the same object for what appears to
be the same reason.

It is worth being explicit that this is not a third position between Reichenbach
and Malament, because the two are not rivals and the substrate says why. They answer different
questions, and both answers stand as given. Reichenbach's question is whether distant simultaneity is
a fact independent of stipulation; the answer is no, and that is the substrate's absent direct
read --- no observation selects $\varepsilon$ because every observation is funded by the throughput
it would have to measure. Malament's question is whether a unique nontrivial simultaneity relation is
\emph{definable} from causal connectibility together with an inertial worldline; the answer is yes,
and that is the substrate's closure cycle stated in causal terms, since causal connectibility is the
round trip. Definable is not measurable, and unmeasurable is not arbitrary. The two results are the
negative and positive halves of one situation --- the direct read is absent, the closure is
present --- which is why both had to come out true, and the appearance of conflict comes only from
reading them as competing answers to a single question.

A third question is then left over, and it is the one the preceding paragraphs answer: not whether
$\varepsilon=\tfrac12$ is factual, nor whether it is definable, but why it is universally
\emph{adopted} rather than merely available. That is the well-formedness answer --- shortest
description, cheapest to maintain --- and it is an account of an inference rather than a further
position on the metaphysics. This deflation is not the consensus reading, and should not be presented
as one: Malament is frequently cited as telling against conventionality rather than as orthogonal to
it, and the disagreement about how much conventionality \emph{matters} survives everything said here.
What does not survive is the impression that the two results contradict.
\end{remark}

\begin{remark}[One move: removing time as a primitive]\label{rem:one-move}
It is worth naming the general move this is an instance of, because the result looks specific and is
not. \emph{Removing time as a primitive is what makes the pair coherent.} Where time is a background
against which events are laid out, distant simultaneity is a fact about the world that a frame may
fail to reach, and the conventionality result reads as an epistemic limitation --- something is
there, and we cannot get at it. Where time is inferred, built by a frame from represented succession and often modelled by
closure cycles rather than supplied as a substrate dimension (Graphic Equalisation's
Section~\ref{GE-sec:time-not-primitive}, \emph{Time is not a dimension of the substrate}), there is
no such fact to fail to reach:
simultaneity is a relation within a constructed axis, and two separated frames have not constructed
one axis between them. Reichenbach's negative is then not a discovery about a hidden quantity but the expected shape of a question that was never well
posed, and Malament's positive is the report of what structure is genuinely there. The pair is what a
substrate without primitive time predicts, rather than something it has to accommodate.

The family is worth seeing together, since its members are usually met one at a time and look
unrelated: a frame's temporal extent is bounded by its own inference rather than by the world; causal
order is prior to time rather than read off it; one-dimensional time is not forced; a block reading
is not prohibited; a frame running no rules has no time at all while things continue to happen to it
from outside; and now, distant simultaneity has no fact behind it to be ignorant of. Each is argued
in its own place. They are one move.

Sized honestly, because the parts of this that are old are most of it. That time is not primitive is
a long relational tradition and this framework is a late entrant to it \citep{rovelli1996}, and the
link between conventionality of simultaneity and relationalism about time is not new either. What is
narrower here is that a single mechanism delivers both halves --- the absent direct read and the
available closure --- so the two results need no reconciling and arrive together. It also cuts as
evidence rather than only as interpretation: a framework holding time primitive must treat the
conventionality as a brute epistemic limit, where one that does not, predicts it. The falsifier is
correspondingly sharp, and it is sharper for being narrower. Exhibit a determination
of the one-way speed in \emph{this} universe that rests on no synchrony stipulation, and what fails is
the empirical claim that our rules supply no unlocked counter --- the cycle unit itself would survive,
since a second dimension's step count is a different unit and not a finer reading of the same one.
\end{remark}

\begin{remark}[Energy is the price of forced concurrency]\label{rem:forced-concurrency}
The same $c$ bounds \emph{concurrency}, which gives the high-energy regime a
reading. Within a frame, change is self-change and strictly ordered, so forcing
two changes onto one rule-step is the paid exception; a particle accelerator is
the machine that pays it, spending energy to push particles' interactions to the
same step, where the frame --- built to sequence --- cannot order them and the
interaction is \emph{decoherent} in that frame. Two changes cannot be crowded
closer than a step, and forcing them there costs more the closer one gets: the
diverging energy above is this same ceiling read at the change-rate. The strong
reading is ``energy is the price of forced concurrency, and forced concurrency is
frame decoherence''; anything with a number is weak.
\end{remark}

\begin{remark}[Two uses of ``decoherence'', and what relates them]\label{rem:two-decoherences}
This paper uses the word in two ways and they should be separated before either does more work. The
sentence above is the \emph{substrate} sense: a frame cannot hold the configuration, so nothing
continues in it as that configuration --- the sense Graphic Equalisation treats generally
(its Proposition~\ref{GE-prop:decoherence-one}), where what varies between cases is the frame and not
the failure. Everywhere
else below it is the \emph{standard} sense: environmental coupling averaging over relative phase so
that branches stop interfering and a preferred basis is selected \citep{zurek2003}. Nothing here
derives that mechanism; einselection is imported, not recovered, and the reader should not take its
appearance as a result of this account.

What relates them is a shape rather than an identity, and stating it that way is deliberate. In both,
a frame fails to hold something as one coherent structure and the difference passes to the container:
the environment absorbs the phase in one case, the ordering the frame cannot supply is available only
above it in the other. What differs is the \emph{outcome} --- the standard case selects exactly one
branch, the substrate case realises none --- and on this account that difference is not in the failure
but in whether the container can absorb what it receives, which is Graphic Equalisation's frame-indexing
applied here. The two are therefore plausibly one operation at different containers, and that is
offered as a reading rather than a derivation: it would be established by showing the selection
follows from absorption alone, which is not shown here.
\end{remark}

\section{Relativistic clock laws and the budget reading}\label{sec:sr-gr}
Special relativity is the flat-spacetime limit of general relativity. The budget picture is a proposed re-description of their clock laws, and requires an explicit map to the metric before it predicts a rate.

\begin{remark}[The budget picture does not derive a clock law]\label{prop:sr-gr}
In flat spacetime the imported relation is $d\tau/dt=\sqrt{1-v^2/c^2}$.
For stationary clocks in a static metric, using signature $(-,+,+,+)$,
$d\tau/dt=\sqrt{-g_{00}}$. Gravitational clock comparison is set by the
metric, not local mass density: a clock inside a hollow massive shell can be
redshifted relative to infinity despite zero local density. Calling mass
\emph{change-load} does not construct $g_{00}$ or recover either formula.
This is the boundary stated in the reconciliation companion's
Proposition~\ref{HG-prop:no-clock-law}.
\end{remark}
The equivalence principle is an additional recovery obligation. Naming inertial
and gravitational mass as one change-load supplies no composition-dependent
prediction or equality proof. The passive-continuation companion states the
required matter-coupling conditions in its Section~\ref{PC-sec:gr}.

\begin{remark}[Unit normalisation does not establish relativity]\label{cor:reconcile}
Setting the numerical value of $c$ to one fixes units. Recovering observed
Lorentz invariance, gravitational redshift and universal matter coupling
requires physical dynamics in addition to that convention. The budget reading
therefore supplies an interpretation target, not a reconciliation theorem.
\end{remark}

\subsection{Relativity and collapse}\label{sec:collapse}
Relativity also reconciles with the collapse picture of quantum measurement, by
the division of Section~\ref{sec:c-rate}: the
continuous in the maps between frames, the discrete in a frame's own states. The continuous
geometry of general relativity is a cross-frame \emph{map}; the collapse of
quantum measurement is the in-frame \emph{event}; and $c$ --- the change-to-time
rate --- is the seam joining the discrete event to the continuous metric. The two
are the map and the event of one change-substrate, not rival ontologies, and the
standing incompatibility is the artefact of treating them as rivals. This is a
claim about the \emph{shape} of the reconciliation, not a derivation: that the
field equations follow is not asserted, and no theory of quantum gravity is
offered.

\section{The measurement problem is frame-relative collapse}\label{sec:measurement}
Collapse in Graphic Equalisation is a local write that realises one continuation
\emph{relative to the frame that interacts}. It is not a physical process needing
a mechanism --- no dynamical collapse, no spontaneous localisation --- but the
in-frame realisation that occurs when frames couple. So there is no
frame-independent ``when'': each frame realises its own outcome upon interaction,
and a frame that has not interacted holds no collapse. This is Axiom~$-1$ and
relational quantum mechanics \citep{rovelli1996} --- facts indexed to
interactions, not assigned from nowhere.

\begin{proposition}[Wigner's friend is the map and the event of two frames]\label{prop:wigner}
The friend measures the system and realises an outcome --- collapse relative to
the friend, who interacted. Wigner, not having interacted with friend-and-system,
holds relative to himself the continuous unitary \emph{map} (Section~\ref{sec:collapse}).
Both are correct: the friend holds the in-frame \emph{event}, Wigner the
cross-frame \emph{map}. The apparent contradiction requires
a single frame-independent fact about whether collapse occurred, which
Axiom~$-1$ forbids; without that demand the friend's definite result and Wigner's
superposition coexist without strain.
\end{proposition}

That there is no observer-independent fact of the matter is, in this literature, a
theorem and not only a stance. \citet{brukner2018} shows that the existence of
observer-independent facts, together with locality and freedom of choice, is
inconsistent with quantum predictions in an extended Wigner's-friend setting, and
\citet{frauchiger2018} show that agents who use quantum theory to reason about other
agents who use quantum theory can be brought to contradictory conclusions about the
same run. Neither result is evidence for the change-substrate --- they constrain
every account equally --- but they remove the option Proposition~\ref{prop:wigner}
declines to take, that of a single frame-independent record the friend's and Wigner's
descriptions must both be answerable to. Axiom~$-1$ forbids that record on structural
grounds; these results forbid it on quantum-mechanical ones, given the auxiliary assumptions each names.

\begin{corollary}[Any coupling is a measurement; outcomes are definite; no special observer]\label{cor:observer}
Any coupling is a measurement relative to the coupled frames, so no conscious or
privileged observer is required and none is distinguished --- a plate, a rock, or
a friend measure relative to whatever they couple to. Outcomes are definite
because collapse realises \emph{one} continuation, and the basis in which they are
definite is set by the distinctions the interaction makes --- the aperture of the
coupling, which is decoherence selecting the pointer basis. There is no cut and
no collapse mechanism to locate.
\end{corollary}

\begin{remark}[Measuring is not observing]\label{rem:measure-not-observe}
\emph{Measure} here is the
coupling itself and not an act: \emph{Is} holds (its Proposition~\ref{IS-prop:one-way}) that what
stands in relation without transformation or grammar can be observed and cannot observe,
and that is consistent with this, because observing in its sense is \emph{making}
a distinction and requires rules, while measuring in this sense is standing in the
relation that makes an outcome definite relative to a partner. A rock does the
second and not the first.
\end{remark}

This dissolves the measurement \emph{problem} --- the when, the how, the who, and
the definiteness. That outcomes are realised frame-relatively is the shape; how
their probabilities are weighted is a separate question, taken up in
Section~\ref{sec:amplitudes}.

\section{Entanglement is dimensional access through frames}\label{sec:entanglement}
An entangled pair is one system in the superframe, described by a joint dimension
--- the correlation --- that is a first-class quantity of the superframe and
\textsc{null} in each subframe individually (Graphic Equalisation's
Proposition~\ref{GE-prop:higher-frame}, higher-frame first-class structure). Each subframe has dimensional access to it only through
its own boundary, and only in part; the full dimension lives one level up, along
the structural-depth axis rather than the spatial one.

That the subframe/superframe split is itself frame-relative is not only assumed here.
\citet{aliahmad2022} prove, in the quantum-reference-frame setting, that the
decomposition of a composite system into subsystems is frame-dependent --- what counts
as ``the parts'' changes with the frame relative to which the whole is described. That result takes frame-covariance as
its premise and derives the relativity of subsystems from it, so it does not establish the
non-privilege of frames --- that remains this framework's own commitment, held as an axiom and
not borrowed as a theorem. What it does establish is that the subframe carving inherits the
non-privilege: no single decomposition into parts is preferred, so the joint dimension below is
a claim about a decomposition, made from inside it, and not a claim about a preferred
carving of the world.

\begin{proposition}[Reduced states of a pure entangled pair]\label{prop:mixed}
For a pure entangled bipartite state, either reduced state is mixed: in its
Schmidt decomposition at least two coefficients are nonzero, so tracing out
the other subsystem yields a density matrix of rank greater than one. This
is a result of the imported quantum state space. In frame vocabulary the joint
pure description is available at the superframe, not locally. A mixed reduced
state alone does not establish entanglement, since a separable mixture can
also have mixed marginals.
\end{proposition}

\begin{corollary}[Entanglement requires the joint quantum state]\label{cor:nonlocal}
For an entangled pair, the superframe represents the joint quantum state.
It is not a classical common-cause variable accessed independently by the two
parties. Under measurement independence, a Bell-local model with factorisation
$p(a,b\mid x,y,\lambda)=p(a\mid x,\lambda)p(b\mid y,\lambda)$ satisfies
Bell inequalities regardless of whether $\lambda$ is called local or
superframe-level \citep{bell1964}. The quantum joint state violates that
factorisation while its local marginals permit no signalling. This is a
re-description of the imported quantum formalism; the substrate vocabulary
alone derives neither the violation nor its magnitude.
\end{corollary}

Quantum mechanics is the standing demonstration that the in-frame description is not exhaustive, and
it is an empirical one rather than an interpretive preference. Correlations violating Bell's
inequality are not functions of the separated subframes' own states \citep{bell1964}, which is
precisely what this paper means by a dimension held at the superframe, and they have been exhibited
without the loopholes that once allowed the conclusion to be resisted \citep{hensen2015}. So
structure beyond a subframe's own description is not merely coherent in this framework; it is
measured. The superframe account is nevertheless a \emph{shape claim}, not an additional measured
variable: it earns the mapping from the measured failure of subframe-local descriptions to a joint
first-class dimension. It fails if a complete account of the observed correlations can be given from
the separated subframes' own state alone, without joint structure or an equivalent containing
relation. Bell tests do not select this vocabulary over every interpretation; they expose the
subframe-local shape that this vocabulary says is unavailable.

\begin{remark}[Ascent yields a different somewhere, not a view from nowhere]\label{rem:ascent-different-somewhere}
The limit on that optimism should be stated in the same breath, because the obvious route is closed
by an established result rather than an open question. No-signalling means entanglement cannot carry
information between separated locations, so it cannot be used to synchronise their clocks, and the
superframe access we demonstrably have is therefore not access to distant simultaneity
(Proposition~\ref{prop:closure-three}). What a superframe does deliver is
worth being exact about, since it is more than nothing and less than an outside: holding both
locations, it has units that close across them, so quantities the subframes cannot obtain are
ordinary in \emph{its} terms --- and being a frame, it is blind to its own runner in precisely the way
they were to theirs. Each level dissolves the one below's restriction and reproduces the restriction
at its own. So ``impossible'' was the wrong word twice over: wrong for a restriction that holds of one
frame and is lifted by the frame above it, and wrong again for suggesting that some sufficient ascent
would lift it altogether, when what ascent yields is a different somewhere rather than a view from
nowhere.
\end{remark}

\begin{remark}[A wiring analogy cannot establish Bell violation]\label{rem:light-switch}
Shared classical wiring is a local common-cause model and does not account
for Bell-violating correlations. A joint quantum state and its measurement
rules supply the distinction required here.
\end{remark}

Honest boundary: seeing the connecting dimension dissolves the \emph{spookiness},
not the \emph{numbers} --- the Tsirelson bound, the correlation strength and the
Born weights are not derived; they become ordinary physics questions rather than
spooky ones.

\section{Amplitudes are a rule on Narrative; recovering the Born rule}\label{sec:amplitudes}
An amplitude is not a foreign number-structure but the directional Narrative of
the momentum paper, acted on by a rule. That paper's directional
momentum $p=\phi_W(\mathcal N)\,\hat{\mathbf d}$ (its
Definition~\ref{NM-def:narrative-momentum}) is a magnitude $\phi(\mathcal N)$
--- the operative-history measure --- times a direction $\hat{\mathbf d}$. The complex amplitude is
therefore the directional Narrative, its dynamics a rule applied to it, and its
phase the direction $\hat{\mathbf d}$ read as position in the cycle
(Section~\ref{sec:c-rate}). This motivates the interpretation of complex amplitudes; the complex state space remains imported until the mapping is constructed.

\begin{remark}[A conjugate-pair interpretation of the Born square]\label{prop:born}
For complex amplitude $a$, $a^*a=|a|^2$ is a real nonnegative quantity.
A closure loop can motivate that pairing as an interpretation of an imported
Born weight. Closure alone does not prove that probability equals this quantity,
or rule out another phase-invariant function of $|a|$. Recovering the rule
requires the additional probability and state-space assumptions below.
\end{remark}
The recovery below uses complex linear amplitudes, phase-averaged probability
additivity and a noncontextual measure on projectors as separate assumptions.
Their substrate mapping remains unproved. In particular, a reduced-density
matrix probability calculation would already use the trace rule. The
assumption issue is explicit in envariance work and its analysis
\citep{schlosshauerfine2005}; a different reconstruction of measurement postulates
retains unitary quantum structure, composition and finite parametrisation
\citep{masanes2019}.



This is a candidate interpretation, not a derivation of the probability rule.

\subsection{An attempt to close it}\label{sec:born-attempt}
The candidate can be pushed further than a candidate, though not all the way. Take
the amplitude as a directed magnitude --- a Narrative's size $r$ times a phase ---
and follow four steps, marking what each earns. \emph{This subsection is an attempt,
flagged as such, not a claimed result.}

\emph{Complex (motivated).} The phase is position on the realisation cycle, and a
cycle's traversal parameter is one-dimensional: $S^1=U(1)$. So an amplitude is
$r\,e^{i\theta}$, a complex number --- real would be a phase-less (cycle-less)
magnitude, and the one-dimensional closure loop gives a single angle, not the three
of a quaternionic phase. Motivated by closure, not forced.

\emph{Linear (a candidate mechanism, not a proof).} Independent contributions to one
outcome superpose by addition. This has a candidate mechanism rather than a bare posit.
Free evolution is the amplitude's own Narrative, and a Narrative is what a subgraph runs
when \emph{no rule fires} --- the free continuation simply carried forward
(the momentum paper's free continuation, its
Definition~\ref{NM-def:free-continuation}, and the free-continuation default of
Section~\ref{sec:noether}). With no rule firing on them, nothing \emph{couples} the components of
a superposition; each free-continues independently, and uncoupled independent continuation
\emph{is} linear superposition, since a nonlinearity would need a rule reading one
component and acting on another --- exactly what does not fire here. Linearity is on this
reading the signature of free continuation. What stays a residual is the rigorous step
from ``uncoupled'' to full additivity and norm-preservation (unitarity), and the
linearity of the imposed supergraph map is not handed over for free: the mechanism
motivates linearity, it does not yet prove it.

\emph{Phase-independent (from no privilege).} A cyclic direction has no privileged
origin, so a global phase rotation is a change of nothing (Graphic Equalisation's
no-privilege principle applied to the phase). The realised probability must
therefore be invariant under a global $U(1)$, hence a function of the modulus alone,
$P=f(|a|)$. This step is earned, and it is the no-privilege thesis --- Graphic
Equalisation's Axiom~$-1$, carried through the mind companion --- doing physics
work.

\emph{Squared (forced, given the above).} For this conditional recovery, assume phase-averaged additivity for alternatives represented as complex amplitude sums. This is an additional probability postulate; it does not follow from decoherence without a probability rule. Averaging the combined modulus over the
relative phase,
\[
\big\langle f(|a\,e^{i\theta}+b|)\big\rangle_\theta = f(|a|)+f(|b|),
\]
since the cross term $2\,\mathrm{Re}(a^\ast b\,e^{i\theta})$ averages to zero. The
square, $f(x)\propto x^{2}$, satisfies it, and the obvious alternatives do not: a quartic
overshoots by the surviving $\langle\cos^{2}\rangle=\tfrac12$ term, the linear
undershoots ($a=b=1$ gives $4/\pi\neq2$), the square lands exactly; normalisation
fixes the constant. This selects the square among the natural candidates rather than proving
it unique --- the uniqueness is what Gleason's theorem supplies below, over all measurement
contexts. The weight is real because closure (a returned cycle) is net-zero phase.

\emph{The Gleason residual, reconsidered.} The internal Narrative is relative in the
\emph{reading}: it is the same standing structure \emph{read as} a vector in the basis the
frame's current rule set fixes, so a superposition is rule-set-relative --- definite in one rule set, superposed in
another --- which is the same \emph{form} as relational quantum mechanics \citep{rovelli1996} without being the
same claim: that programme relativises a physical fact to an interacting physical system, never to a
descriptive framework or rule set, so what is borrowed here is the shape of the relativisation and not
its warrant. Noncontextuality means that the probability assigned to a projector does
not depend on the orthogonal decomposition containing it. Equal status of
measurement bases does not establish that condition. Given a Hilbert space
and a normalised additive noncontextual measure, Gleason's theorem supplies
the trace-rule uniqueness for dimension at least three \citep{gleason1957}.
Those are imported mathematical premises, not consequences of the
no-privileged-frame discipline.

\emph{The Hilbert posit.} The cross-frame reading motivates an inner product,
but does not construct one. Recoverability is injectivity on the operative state
space; it supplies neither linearity nor preservation of norms and inner products.
Likewise a continuous phase parameter does not prove completeness. Complex Hilbert
space, a specified composition rule, and unitary evolution remain premises of this
conditional grounding.

The operational reconstructions \citep{hardy2001,chiribella2011,masanes2011}
can establish Hilbert structure and Born probabilities once their full hypotheses
hold. The Born note's Section~\ref{BN-sec:axioms} records the proposed mapping.
Hardy's continuity requirement connects arbitrary distinct pure states by continuous
reversible transformations; rotating a ray's global phase does not meet it.
CDP purification requires both existence and essential uniqueness under reversible
transformations on a purifying system. Subframe nesting alone proves neither.
Local tomography and complex composition also need separate constructions.
These are proof obligations, not native consequences already supplied.

\subsection{A common interpretation and its limits}\label{sec:born-sr-gr}
Born probabilities and relativistic clocks can be discussed in the same substrate
vocabulary. That common vocabulary establishes no dynamical unification.
Remarks~\ref{prop:sr-gr} and~\ref{cor:reconcile} record the metric assumptions
missing from a change-budget interpretation; the Born mapping records different,
quantum assumptions. No reconciliation of quantum theory and gravity is proved.

\section{Interference and the double slit}\label{sec:interference}
The double slit is obvious once two established facts are held together: causality
is ordering \emph{in a frame}, and a superframe applies changes to a subframe
without passing through the subframe's admissibility.

A superframe \emph{applying} a change to a subframe is not itself a measurement
coupling: it makes no distinction the subframe holds, so it realises nothing for
that subframe --- which is why the superframe forcing below leaves an interference
pattern intact while a which-path coupling destroys it.

\begin{proposition}[Superposition is the possibility structure; interference is amplitudes combining by phase]\label{prop:interference}
Superposition is Graphic Equalisation's possibility structure $\mathcal P(G)$ ---
the admissible continuations a graph holds before collapse --- so a particle ``in
superposition'' is a graph whose $\mathcal P$ has more than one admissible
continuation, none yet realised. Each continuation carries a directional Narrative
(Section~\ref{sec:amplitudes}), an amplitude with a cyclic phase, and before
collapse the alternatives' amplitudes combine, constructive where the phases align
and destructive where opposed. The outcome weight is the combined closed loop,
$|\sum\text{amplitude}|^{2}$.
\end{proposition}

Two established facts make this un-mysterious. First, causality is ordering
\emph{in a frame}: there is no global fact about ``which path the particle really
took,'' because which-path is \emph{undetermined}, not \emph{undisclosed}, until a
coupling realises it --- a question with no answer of any kind until a coupling settles
it. The ``both paths'' is not a particle in two places but the frame's
$\mathcal P$. Second, a superframe applies changes to a subframe \emph{without}
passing through the subframe's admissibility (Graphic Equalisation's
Proposition~\ref{GE-prop:self-governance}: admissibility is self-governance, not armour --- a graph's
grammar governs its self-originated transformation, not world-caused forcing). So the interference is the superframe
applying its combined-amplitude structure to the particle; the particle does not
select a path by its own rules, and cannot veto the combination.

\begin{corollary}[The double slit, duality, and delayed choice]\label{cor:doubleslit}
Both paths are admissible continuations in $\mathcal P$; the superframe applies
their combined amplitudes at the screen, giving fringes. A which-path coupling
collapses before the screen --- one continuation realised, the other no longer in
the sum --- so the superframe applies a single-path change and the fringes vanish.
Wave--particle duality is the map/event distinction: the ``wave'' is the amplitudes
combining (the map, the held $\mathcal P$), the ``particle'' the realised outcome
(the event, one continuation collapsed), and ``which is it really'' names no
difference. Delayed choice adds no retrocausality: collapse attaches to the actual
interaction, not a retroactive rewrite (cost is a local write; no write reaches
the past), so interference-or-not depends only on whether the which-path coupling
occurs.
\end{corollary}

Honest boundary: this grounds the shapes --- superposition as possibility,
interference as phase-combination, the double slit, duality, and delayed choice
without retrocausality. The quantitative fringe pattern needs the amplitude phases
computed, i.e.\ the Hilbert formalism, which is the Born-rule frontier of
Section~\ref{sec:amplitudes}.

\section{Dimensional constants are frame-units; dimensionless ones are not}\label{sec:constants}
The reading of $c$ as a frame's own base unit generalises. Each dimensional
constant is the base unit of one of a frame's own economies --- $1$ in-frame, a
number only cross-frame.
\begin{itemize}[leftmargin=*]
\item $c$ is the change-rate unit (Section~\ref{sec:c-rate}).
\item $\hbar$ sets the scale of the canonical commutator
$[x,p]=i\hbar$, which implies $\Delta x\,\Delta p\ge\hbar/2$
for state preparation. Interpreting this as a finite-aperture budget requires
an additional model. Measurement error and disturbance obey separately stated
relations and are not identical to preparation uncertainty.
\item $k_{\mathrm B}$ is the information--energy conversion: Landauer's
$k_{\mathrm B}T\ln 2$ \citep{landauer1961} --- the thermal lower bound for erasing one unbiased logical bit. Bijective updates need not pay this erasure floor; physical costs require their own accounting.
\item $G$ is the change-load--geometry coupling, mass deforming the containing
frame.
\end{itemize}
That physics sets $c=\hbar=k_{\mathrm B}=G=1$ (natural units) is then not a
convenience but a recognition: each is a frame-unit, read as $1$ in the frame's
own terms with its familiar magnitude a cross-frame artefact
(Proposition~\ref{prop:c-unit}).

\begin{proposition}[The dimensional/dimensionless boundary]\label{prop:dim-boundary}
A dimensional constant is a frame-unit and reads as $1$ in its own terms; a
dimensionless one --- the fine-structure constant $\alpha\approx1/137$, the mass
ratios --- cannot be set to $1$ by any unit choice, being a genuine cross-frame
\emph{ratio}. The framework says nothing about their values. About their type it
says: a constant is a ratio or it is a unit artefact, and there is no third kind.
There are no magic numbers --- standalone quantities that simply are, attached to
nothing.
\end{proposition}

The type claim is forced rather than preferred, and the forcing is at the floor. A ratio is a
relation between two quantities: two ends and something carried across. A magic number is a
\emph{label} --- a value borne by one thing, answering to nothing else --- and the base has no labels
for it to be written as, carrying relata and relations and no marks on either (\emph{Is}'s floor,
made explicit in its Section~\ref{IS-sec:graph}). \emph{So a standalone constant has nowhere to
live.} Where a theory carries one anyway, it is carrying a value whose relata it has not identified, which is a write nobody can read
back --- the silent fourth cost of Graphic Equalisation's well-formedness
(Definition~\ref{GE-def:well-formed}) --- and the defect is in the
theory rather than in the world.

What this does to the open question is change its shape and not its status. \emph{It explains no
constant's value and does not pretend to.} What it denies is that the value is the right thing to ask
after first. ``Why is $\alpha$ about $1/137$'' asks for the provenance of a number; on this account
the number is what a relation carries, so the question is which two things stand in it and what fixes
their relative size --- and the standard framing, in which the Standard Model has some two dozen free
parameters, is then not a list of two dozen mysteries but a list of two dozen relations whose relata
are unnamed. That is a harder question and a better-posed one, and nothing here answers it.

The commitment is checkable in the one direction that matters. \emph{Exhibit a dimensionless constant
that is a ratio of nothing} --- irreducibly standalone, not expressible as one quantity against
another even in principle --- and the type claim fails. The candidates all survive it so far, and
survive it trivially: $\alpha$ compares a coupling against $\hbar c$, the mass ratios are ratios by
construction, $\theta_{\mathrm{QCD}}$ is an angle and an angle is an arc against a radius, $\Lambda$
in Planck units is one energy density against another. Counts are the case worth naming, since a
count looks like the counterexample: the number of spatial dimensions, or of generations. A count is a ratio to the unit of what is counted, which is degenerate but is not standalone --- there is still
something it is a count \emph{of}, and that is the second end.

\emph{And the consequence that makes this empirical rather than taxonomic: a constant is constant
except when the things that set it change.} Constancy is a property of the relation's stability and
not of the number, so \emph{constant} names a state that can fail rather than a category that
cannot. Nothing in this account enforces it. That is not a prediction that constants do vary --- it
is the removal of a reason to expect they cannot, and exact constancy maintained forever would be a
fact wanting an explanation rather than the default it is usually taken for.

The split above then does work it is not usually asked to do. \emph{Only a ratio can vary, because
only a ratio is anything.} A claim that $c$ has changed is not a claim about the world: $c$ is the
frame's own unit, read as $1$ in its own terms, and a varying ruler measured by itself is no
measurement --- which is \citeauthor{duff2002}'s point arrived at from this side. A claim that
$\alpha$ has changed is a claim with content, because $\alpha$ compares two things and one of them
could have moved. So the varying-constants question is well posed exactly for the dimensionless
remainder and ill posed everywhere else, and this is why searches quote $\Delta\alpha/\alpha$ and
never $\Delta c$.

\emph{What a measured variation would be is a measurement of the relata}, which is the part worth
having. If $\alpha$ drifts, the relation between the quantities it compares has changed, and the
structure of the drift --- with position, with epoch, with environment --- constrains which of them
moved. On this account that is not an embarrassment to be explained away but the one observation
that would name the ends of a relation the theory currently carries unnamed. The empirical position
is that the bounds are tight: reactor, clock-comparison and nucleosynthesis constraints all limit
any drift severely \citep{uzan2011}, and the reported spatial dipole in $\alpha$ from quasar
absorption \citep{webb2011} is contested and not independently confirmed, so nothing here rests on
it. \emph{It is named as the shape a confirmation would have}, not as evidence.

\emph{And a ratio can look constant without being constant, which the tightness of those bounds does
not distinguish.} Write the ratio $R=A/B$. Its fractional response is
$\delta R/R=\delta A/A-\delta B/B$, so a large denominator does not suppress a fractional change in
$B$ --- the condition worth stating rather than hiding. What a large denominator suppresses is an
\emph{absolute} one: a perturbation of fixed size $\delta$ enters as $\delta/B$ and vanishes as $B$
grows. Wherever what can move a quantity is bounded in absolute terms, which is the ordinary case, a
ratio with a large enough denominator reads as constant whether or not it is.

That is the finite aperture of $\hbar$ above, arriving at the constants. A frame cannot resolve a
change below its own cell, so \emph{constant} and \emph{varying beneath resolution} are the same
reading to it, and constancy joins the list of properties that are aperture-relative rather than
intrinsic. What is constant to one frame need not be to a finer one, and no frame is entitled to read
its own resolution as a fact about the world (Axiom~$-1$ again, at the measurement).

Two consequences follow and the second is the useful one. First, the empirical bounds say exactly
what they say and not more: they bound variation \emph{below the resolution achieved}, which is a
statement about the aperture of the measurement and not about the world. Nothing here disputes them
and nothing licenses reading them as an absence.

Second --- and this is where the diagnostic hope of the previous paragraph becomes a search strategy
rather than a wish --- \emph{the constants most likely to yield a measurable variation are not the
ones deemed most fundamental, but the ones whose denominators are smallest against the perturbations
available to them}. Apparent constancy is then predicted to track denominator size rather than
depth, and a programme aimed at whichever constant seems most basic is aimed by the wrong criterion.
This account cannot rank the candidates, having no access to the relata that would fix the
denominators; what it supplies is the ordering principle and the claim that it is the right one.

One consistency check on the section before it. The decay constant of
Section~\ref{sec:decay} is a rate, and a rate carries a dimension, so by this proposition it is a
frame-unit and not a fact about the world either. What would be a fact is the ratio of one decay rate
to another, and it is that ratio which the argument there requires to be stable --- not the
half-life in anybody's seconds.
It sharpens the real question rather than answering it: the numbers that need
explaining are the dimensionless ones, while the pseudo-mysteries --- why $c$,
why $\hbar$ --- are unit conventions. This is the standard observation that the
number of fundamental \emph{dimensional} constants is a matter of convention
\citep{duff2002trialogue}, landing on the side that they are units; the dimensionless
remainder is left exactly as open as it was found.

\section{The Planck scale, quantization, and fine-tuning}\label{sec:planck}
Three consequences follow. First, the \emph{Planck scale} is where the
\emph{cross-frame ratios} between the four economies all reach unity together
($\ell_P=\sqrt{\hbar G/c^{3}}$ and its companions). Each constant reads $1$ in its own
frame always, so this is not a scale at which they become $1$; it is the coincidence,
visible only in the comparison between economies, of the aperture cell ($\hbar$), the
change-rate ceiling ($c$) and the change-load--geometry coupling ($G$) --- an artefact of
comparison in the same sense the SI value of $c$ is. Below it a
frame cannot resolve, cannot order, and horizons its own geometry, so nothing
there is representable --- the ``\,?\,'' of Graphic Equalisation's epistemic
boundary with a location. The Planck scale \emph{is} the frontier, which is why
physics predicts its own breakdown there: not a failure of theories but the
exhaustion of a frame's units, below which the pre-relational substrate begins.

Second, $\hbar$ sets the action scale of canonical commutators. Discrete
energy spectra require the system's Hamiltonian and boundary conditions; a
free particle has a continuous spectrum. Neither energy nor action generally
comes in universal $\hbar$-steps.

Third, the \emph{dimensionless} constants are contingent rule-content: the
framework cannot derive $\alpha$ or the mass ratios, but it can say what kind of
thing they are --- the contingent specifics of this substrate's rule-set (Graphic
Equalisation's Remark~\ref{GE-rem:nothing-compulsory}), so ``why $\alpha$'' is not ``why this necessary
value'' but ``why this rule-set,'' answered by self-location
(Proposition~\ref{prop:self-location}). Two type claims are now in play and they agree: a
dimensionless constant is contingent rule-content, and it is a \emph{ratio} rather than a label
(Proposition~\ref{prop:dim-boundary}).

The second sharpens what a fine-tuning argument is doing. \emph{Such an argument treats a constant as
a dial} --- a value that could have been set otherwise with everything else held fixed --- which is
to treat it as a magic number, the one thing this account says no constant is. Turning a dial on a
ratio is not holding everything else fixed: it is altering the relation between the two quantities
the ratio compares, and whether that is a coherent counterfactual at all is the question such
arguments skip rather than answer. That is a second reason the framing misprices its question, and it
is independent of the first --- self-location would dissolve fine-tuning even if constants were
labels, and this would bite even if self-location failed. Both relocate fine-tuning at its root,
which Section~\ref{sec:cosmology} takes up.

\section{Symmetry and conservation}\label{sec:noether}
A symmetry is a transformation that produces no residual: a change the structure
does not feel, free and reversible, a no-op in the container's grammar. Noether's
theorem, symmetry $\leftrightarrow$ conserved quantity, is then the cost side of
that. A quantity is conserved precisely when altering it would cost --- produce
residual, require a paid write --- so the symmetric direction is free and the
quantity perpendicular to it cannot move without paying. What a transformation
leaves untouched is what it costs to change.

\begin{proposition}[Momentum conservation is the free-continuation default]\label{prop:momentum-cons}
Space-translation symmetry --- the rules the same everywhere --- produces no
residual under translation, and momentum is conserved. Momentum conservation is
the free-continuation default Graphic Equalisation states (its
Remark~\ref{GE-rem:collapse-write}): the un-rewritten Narrative, Newton's first
law, continuation costing nothing. Noether's momentum case and inertia are one
statement.
\end{proposition}

\begin{proposition}[Energy conservation is contingent on time-translation symmetry]\label{prop:energy-cons}
Time-translation symmetry --- the rules the same each cycle --- conserves the
change-budget, which is energy. The failure case is checkable: where the rules
effectively change across cycles there is no time-translation symmetry and energy
is not globally conserved. This is the standard general-relativistic result that
a time-dependent spacetime has no global energy conservation --- in an expanding
universe photons redshift their energy away --- here from the same principle:
rules changing across cycles (cosmic expansion, or minting) means no conserved
energy.
\end{proposition}
Charge conservation follows from the \emph{global} phase symmetry --- a
uniform re-description carrying no in-frame residual, which is the shape above.
\emph{Local} gauge symmetry is a different thing: a redundancy of description
rather than a symmetry of states, and by Noether's second theorem it yields
identities and constraints rather than an independent conserved charge. The
framework's reading applies to the first; the second is the cross-frame freedom
(in the maps, where the continuous lives) that leaves the physics untouched, and
carries no charge of its own to conserve.

\begin{remark}[The gauge group is contingent rule-content]\label{rem:gauge-contingent}
The framework grounds \emph{that} a symmetry yields a conservation law --- the
shape, residual-free $\leftrightarrow$ cost-to-change --- but \emph{which}
symmetries hold, the $\mathrm{U}(1)\times\mathrm{SU}(2)\times\mathrm{SU}(3)$ of
the Standard Model, is a fact about this substrate's rule-set, not a necessity.
``Why this gauge group'' is ``why this rule-set,'' the same relocation as ``why
$\alpha$'' (Section~\ref{sec:planck}), answered by self-location
(Proposition~\ref{prop:self-location}). The conservation \emph{structure} is
universal; the particular symmetries are this universe's.
\end{remark}

\section{The second law}\label{sec:second-law}
Entropy and its arrow follow from the same economy.

\begin{definition}[Entropy is what a frame does not distinguish]\label{def:entropy}
The entropy a frame holds is the log-count of the microstructure it carries but
cannot tell apart --- what falls below its aperture. Entropy is therefore
frame- and aperture-relative from the start: a finer aperture resolves more and
reads less entropy, a coarser one more.
\end{definition}
This is the Gibbs--Jaynes reading, in which entropy is a property of the
observer's coarse-graining rather than of the world in itself \citep{jaynes1957};
here it is grounded rather than assumed --- entropy is \emph{what this frame does
not distinguish}.

\begin{remark}[Entropy growth requires a statistical grounding]\label{prop:second-law}
Reversible microscopic dynamics preserves fine-grained Gibbs entropy but
can increase coarse-grained entropy, as in free expansion from a restricted
initial macrostate \citep{jaynes1957}. A low-entropy preparation,
coarse-graining and dynamical assumptions supply the statistical arrow.
Logical erasure has a thermodynamic cost; summing Landauer bounds does not
derive the second law. A predecessor ceasing to be represented is neither
necessary nor sufficient for coarse-grained entropy to rise.
\end{remark}

\begin{remark}[Absence is not residual]\label{rem:absence-not-residual}
A clarification is needed here, because Graphic Equalisation separates two things this
could run together. What falls below an aperture is \emph{absent} to that frame rather than
held-and-unresolved, and absence is not residual --- they have opposite consequences there. Entropy as
defined above is the microstructure a frame cannot tell apart, which is absence at its grain, and it is
\emph{not} the same quantity as un-absorbed residual, which is inadmissibility a grammar did see and
left in place. The two rise together under coarsening and remain distinct in kind.
\end{remark}

\begin{remark}[The increase is one-sided accounting]\label{rem:one-sided}
One further separation, and it is what makes the proposition a reading rather than a law. Demotion
below an aperture is not loss: what falls below one frame's grain is above another's, so the transfer
\emph{balances} across the boundary and is unbalanced only when accounted from one side. The
increase is therefore a property of \emph{one-sided accounting}, not of the whole --- a cross-frame
quantity read in-frame, which is the category this programme's own discipline says produces exactly
this kind of artefact. That is why the law is asymptotic and statistical rather than exact, and why it
can be locally reversed by anything that pays: the payer is the other side of a ledger that was
balanced all along.
\end{remark}

\begin{remark}[A relocation of the arrow's source, not a new arrow]\label{rem:arrow-relocation}
Nothing in that is new physics. Open-system thermodynamics already has entropy falling locally while
rising globally, and the demon is already answered by accounting for the demon. What this account adds
is only where the boundary is: the boundary is a \emph{frame}, so the asymmetry is not a fact about
systems but about individuation, and the second law is what the balance looks like from inside one of
the two frames it balances between. The claim is a relocation of the arrow's source, not a new
arrow.
\end{remark}

\begin{corollary}[The demon is exorcised, and local order is affordable]\label{cor:demon-life}
Landauer's $k_{\mathrm B}T\ln 2$ (Section~\ref{sec:constants}) is the price of
\emph{erasing} one bit --- pushing one distinction back down into the
un-distinguished, which is the discarding write of
Proposition~\ref{prop:second-law}. Maxwell's demon is exorcised by exactly this:
a demon with finite memory must erase to keep operating, and that erasure is the
entropy it thought it removed for free \citep{bennett1982}. Acquiring the
distinction need not cost --- a copy into blank structure discards nothing ---
which is why the cost sits in the reset and not in the observing, and why on this
account it is the \emph{no-copy} asymmetry and not the act of distinguishing that
carries the arrow. Yet local entropy reduction is
possible: minting resolves residual and builds structure, but it is a write, so
it is paid, and the payment is Landauer heat exported to the surroundings ---
local order bought by global entropy increase. This is Schr\"odinger's ``life
feeds upon negative entropy'' \citep{schrodinger1944} grounded: an organism, a learner, a
growing graph is a region that pays to mint local structure and dumps the cost as
heat. Life does not violate the second law; it is the second law with a subsidy.
\end{corollary}

\begin{remark}[The arrows unify]\label{rem:arrows}
The thermodynamic arrow (entropy up), the causal arrow (change enables forward),
and the psychological arrow (memory is present structure of a discarded past) are
not three arrows that happen to align --- but the unification is between two of
them, and the third joins for a different reason. The thermodynamic and psychological arrows are one thing: the
accumulation of change under no-copy, memory being present structure of a
discarded past and entropy being the same discarding counted. The \emph{causal}
arrow is not that. It is the asymmetry of enablement itself --- a change makes
another possible and the second does not thereby make the first possible --- and
it holds whether or not a predecessor survives, which is why Graphic Equalisation
rests the event order on it rather than on no-copy. So two arrows are identical
and the third agrees with them without being them, and the ``why do they agree''
puzzle is answered in two parts rather than one: two descriptions of a single
asymmetry, and a second asymmetry that runs the same way because the changes
carrying it are the changes doing the discarding.
\end{remark}

\section{Decay is determinate; the phase is unretained}\label{sec:decay}
Radioactive decay is the standard exhibit for irreducible chance: the law is exact in aggregate and
the individual event unpredictable. On this substrate the two halves come apart, and the interesting
one is the aggregate.

\begin{proposition}[Exponential decay is what an unretained phase looks like]\label{prop:decay}
A decaying nucleus is a frame with its own closure cycle, and its decay is determinate in its own
cycle count. What no frame holds is \emph{how far through that count it is}: a cycle has no fixed
length and each does a variable amount of change (Graphic Equalisation's
Proposition~\ref{GE-prop:cycle-unit}, that the unit of inferred time is a cycle and requires
closure), so there is no frame-free conversion from the
nucleus's count to the observer's. Retention is paid, and nothing here pays to retain the phase.
\emph{So there is no age for a decay probability to depend on}, and a hazard with nothing retained to
depend on is constant \emph{with respect to everything the nucleus holds}.
\end{proposition}

The qualification in that last clause is load-bearing rather than cautious. A hazard constant
\emph{simpliciter} is exponential decay and nothing else, and exact exponentiality turns out to
require something this account denies --- which is the subject of the deviations below.

The qualification is what earns the proposition, because it is where the account says more than its
rival. \emph{Irreducible chance does not predict a distribution.} That decay is undetermined is
consistent with any law whatever --- peaked, heavy-tailed, bimodal --- and the exponential has to be
put in by hand alongside it, as an observed fact about nuclei rather than a consequence of their
being undetermined. Here the exponential is forced: memorylessness is not an extra postulate but the
absence of a retained quantity, and constant hazard has exactly one survival law. \emph{A nucleus
that carried its own age would be a more expensive object}, and Graphic Equalisation's
write-economics (Proposition~\ref{GE-prop:economics}) says who pays for the carrying. A determinate countdown --- a nucleus that knows its remaining ticks --- is the
model this proposition is most likely to be mistaken for, and it is refuted by the data it is meant
to explain: survivors of a countdown are nearer their deadline, the hazard rises with age, and the
observed law would be peaked rather than exponential.

Nothing here is a local hidden variable of the kind Bell forbids. Decay times of separated nuclei
are not a correlation between separated measurements, so no inequality is engaged; and where one
\emph{is} engaged, the account's answer is the superframe dimension of
Section~\ref{sec:entanglement} and not this phase, which must not be promoted into one.

\emph{The exposure is the deviations, and they are real.} Standard quantum mechanics predicts
departures from exponentiality at both ends: quadratic behaviour at short times, and a power-law
tail at long ones \citep{khalfin1957}. Both have been reported --- the short-time departure in the
tunnelling of ultracold atoms \citep{wilkinson1997}, the long-time one in the luminescence decay of
organic molecules \citep{rothe2006}. A constant hazard admits neither, so the proposition owes an
account of both.

It has one for the short end and it is a prediction rather than a repair.
Graphic Equalisation's cycle-unit result (Proposition~\ref{GE-prop:cycle-unit}) makes inferred time a
\emph{cycle} and requires closure for it, so
below one closure cycle of the \emph{observing} frame there is no inferred-time interval for a law
to be exponential in. Exponentiality is a statement in inferred time; short of a cycle there is no
such statement, and a departure there is expected on this account and unexplained on the bare
irreducible-chance one.

The long end inverts, and the inversion is worth stating carefully because it is easy to overstate.
Khalfin's argument is that exactly exponential survival requires the state's energy distribution to
be exactly Lorentzian, and a Lorentzian has support on the whole line, including arbitrarily far
below any threshold; a Hamiltonian with a ground state is bounded below, so the distribution vanishes
on a half-line and its transform cannot decay exponentially \citep{khalfin1957}. \emph{Exact
exponentiality therefore requires a spectrum unbounded below} --- which is to say, it requires a
continuum quantity of unbounded support to be a frame's own state distribution. This paper has
already denied exactly that: continuum structure lives in the maps between frames and never in a
frame's own stepped states (Section~\ref{sec:c-rate}). So exact exponentiality is a cross-frame
idealisation that no frame exhibits, and a departure from it at long times is what this account
should have predicted rather than what threatens it.

\emph{What that earns is the deviation in kind, and not its form.} The power law and its exponent
follow from the threshold branch point, which is a fact about the composite system's spectrum and not
about what the nucleus retains; nothing here derives them, and treating the inversion as though it
did would be claiming the result rather than the direction. So the standing exposure is narrower than
it was and is still an exposure: \emph{if the long-time tail is confirmed with an exponent this
account cannot reach, the proposition is bounded and not exact.} What is no longer available as an
objection is that the deviation exists at all.

One further thing the account does supply is the \emph{visibility} structure, which is worth
separating from the derivation it does not supply. The deviation is not absent through most of a
decay and then present at the end. It is there throughout and sits beneath the aperture
(Section~\ref{sec:constants}), and what changes at long times is not the tail but the exponential: a
power law and a falling exponential cross once $e^{-\Gamma t}$ has dropped below $t^{-n}$, which is
many lifetimes in: \citet{rothe2006} report the crossing at around ten lifetimes. The
exact crossing depends on the tail's prefactor and is not estimated here. \emph{So the tail becomes
visible because the thing hiding it decays away, not because the tail grows}, and ``deviations appear
at long times'' is a fact about what is resolvable rather than about what is happening.

\emph{That is not the missing derivation.} The gap is that the constant hazard here is a claim about
what the \emph{nucleus} retains, while the deviation comes from the \emph{composite} spectrum being
bounded below. Those are different objects. What reconciles them is that the floor is a fact about
\emph{any} frame, so it holds of the composite and of the nucleus alike: the constant hazard is what
the nucleus retains, the required deviation is what the container's floor imposes, and neither is a
claim about the other.

\begin{proposition}[The spectrum is bounded below because absence is]\label{prop:spectrum-floor}
A frame's spectrum is bounded below because absence is: occupation is a count, so it is non-negative,
and the sign of a frequency is an orientation rather than a rate. With both,
$H=\sum_k n_k\hbar\omega_k\ge 0$ term by term. What is derived is that a floor \emph{exists}, not
where it reads.
\end{proposition}

Stepping alone gives no floor: the integers are stepped and run both ways. \emph{One substrate fact is needed and it is applied in two places}, the second of which is where
the sign problem hides.

\emph{Occupation is a count, so it is non-negative.} A frame's energy is carried by quanta present,
$\sum_k n_k\hbar\omega_k$, and each $n_k$ is a cardinality of what is there. Absence is zero and not
a negative quantity; there is no less-than-nothing to count down into. The substrate's unbounded
direction for a count is the extensive frontier --- always one more step outward --- which is
$\mathbb{N}$, floored at zero by construction, never $\mathbb{Z}$. What is required here is that
occupation is a count, not that the ladder is evenly spaced; the $\hbar$-steps of
Section~\ref{sec:planck} are the action quantum and are not needed for this.

\begin{remark}[The sign of a frequency is an orientation, not a rate]\label{rem:sign-is-phase}
\emph{And the sign of a frequency is not a rate.} An $\omega_k$ appears as $e^{\mp i\omega t}$ with
either sign, and this is where the negative-energy problem historically entered. The negative case is
\emph{not} a cycle run backwards: it is the same cycle traversed in the opposite \emph{phase},
running forward like any other. \emph{Negative time does not exist and phase does}, which is the
distinction Graphic Equalisation states in general (its Remark~\ref{GE-rem:phase-not-time}, that a
cycle reads two ways --- how many, and where within): a count of completed cycles has no direction to run backwards in, while a phase is
oriented and carries no duration at all. Reading the second as the first is what manufactures a
backwards time, and the manufactured object is then what needs forbidding. So the sign is an \emph{orientation} --- which way round the relation
is read, carried by the edge and not by any magnitude --- while the rate itself is a count of cycles
against the frame's own cycle (Graphic Equalisation's Proposition~\ref{GE-prop:cycle-unit}),
non-negative for the reason
occupation is. \emph{So this is one premise applied twice and not two}: absence floors the count, and
the apparent second sign was never a quantity. What looked like a negative-energy state is a
positive-energy state read in the opposite orientation --- the reinterpretation physics adopted, with
the sea of filled states as the alternative that is not needed here.
\end{remark}

That is both premises, and the proposition follows.

The zero of energy is a choice of origin,
and an origin is read in-frame exactly as a unit is (Section~\ref{sec:constants}), so a spectrum running to
arbitrarily negative numbers is no counterexample: a binding energy is negative against a
conventional zero and the system is still bounded below by its own ground state, which is the floor
this proposition is about. The claim is boundedness, which is invariant. The number at which the
bound sits is not, and nothing here needs it to be.

That is the \emph{spectrum condition} --- positivity of the energy in every frame, which is where it
is usually taken as an axiom \citep{streater1964} --- arriving as a consequence rather than a
postulate: each frame floors its own spectrum for the same reason, and none is privileged, and it
closes the
argument the other way round from how it was posed. Exact exponential survival needs spectral support
running to $-\infty$; the floor forbids it; so \emph{the deviation is not merely permitted by this
account but required by it}, and required by the same fact about absence the substrate uses
everywhere else. So the leading behaviour is exponential because nothing is retained, and the
departure is required because a floor exists --- two facts about two objects, not one derivation
across them.

\emph{The exponent is a different matter and the account is right not to fix it.} For a spectral
density behaving as $(E-E_{\min})^{\beta}$ just above the floor, the survival amplitude falls as
$t^{-(1+\beta)}$ and the probability as $t^{-2(1+\beta)}$, so the exponent is set by $\beta$ --- how
many distinguishable configurations sit just above absence. That is a fact about how the frame's
dimensions are arranged, which is contingent rule-content (Section~\ref{sec:dimensions}) and not
substrate structure. \emph{So the substrate fixes that there is a tail and cannot fix its exponent,
and the reason is principled rather than a shortfall}: the exponent is contingent, and a substrate
result that delivered a universal one would be claiming more than the substrate knows.

Which converts the exposure into a test. \emph{The exponent should track the density of states at
threshold and not be universal} --- differently structured systems should tail differently, and the
free-particle $s$-wave case with $\beta=\tfrac12$ giving $t^{-3}$ is one value among the available
ones rather than the value. A universal exponent, the same across systems with different threshold
structure, would falsify this.

One identification is doing the work and is worth marking rather than burying: \emph{in-frame energy
is a count of quanta}. That is clean in the occupation-number picture and less clean for interacting
systems, where the count is not sharp and the physical states are dressed. Nothing above survives if
that identification fails for the interacting case, and this section does not establish that it
holds there.

\section{Cosmology: the initial state, the seed, and self-location}\label{sec:cosmology}
By Graphic Equalisation's Proposition~\ref{GE-prop:no-start-end} the event order has no intrinsic
first, so the beginning of the universe is not an intrinsic edge but a
\emph{chosen entry} --- a boundary condition imposed from the containing frame.
This is why physics supplies initial conditions separately from its laws: the
laws are the ordering, which fixes no start, so the entry must come from outside.
The Big Bang is that entry, the seed; what is ``before'' it is not an earlier
time but the frontier, the pre-relational substrate, unrepresentable from within,
the same ``\,?\,'' that sits below the Planck scale (Section~\ref{sec:planck}).

\begin{proposition}[The Past Hypothesis is a seed, hence a hypothesis]\label{prop:past-hyp}
The low-entropy initial state from which the second law takes its direction ---
the Past Hypothesis --- is a seed: the initial rule-and-state configuration
installed from the container. Its long-standing status as an assumption rather
than a derivation is therefore correct in kind, since a seed cannot be derived
from within (a mint cannot authorise its own installation). The arrow of time
then splits cleanly: its \emph{mechanism} is no-copy accumulation
(Section~\ref{sec:second-law}), and its \emph{origin} --- that there is a
low-entropy state to rise from --- is the seed. The direction points away from
the imposed low-entropy start; the running is no-copy.
\end{proposition}

\begin{proposition}[Why this seed: self-location]\label{prop:self-location}
Why \emph{this} seed --- and, equally, why this $\alpha$, this gauge group --- is
dissolved rather than declined, and what dissolves it is self-location. A map is referenced from
``you are here,'' and the seed is the here: the origin the frame is drawn from,
as $c=1$ is the unit it is measured from. The observer is not independent of it
--- by frame-relative identity a different seed is a different observer --- so
``why this seed rather than another'' assumes a trans-seed observer who could
compare, and there is none (Axiom~$-1$). It is that way because we are this way,
and we are this way because it is that way: one determination read from inside,
not two facts one explaining the other.
\end{proposition}
This inverts the fine-tuning inference. The constants are not tuned \emph{to} a
pre-existing observer; the observer is what those constants produce, so the
apparent coincidence is the co-determination of observer and seed mistaken for a
tuning aimed at a target. What remains out of reach is only \emph{derivation}: no formula
computes $\alpha$ or the seed. The \emph{why} is answered; the \emph{value} is
not derived. Reference points take no ``why'' --- they are where the why is
asked.

\section{Dimensions are contingent; time is optional}\label{sec:dimensions}
Time is not a must. A frame whose rules are not firing --- inert, a rock --- has
no time at all: nothing happens in it, there is no succession, and in its own
frame it merely \emph{is}. What an external clock calls ``later'' is, for the
rock, only another \emph{is}; the weathering it undergoes is the \emph{supergraph's}
succession applied to it from above (Graphic Equalisation's
Proposition~\ref{GE-prop:self-governance}: admissibility is self-governance, not armour), not the
rock's own. Time belongs to active, inferring frames --- those
with a firing closure loop to iterate (Graphic Equalisation's Proposition~\ref{GE-prop:cycle-unit})
and the
inference to represent the succession. The rock has neither, so it has no time;
its rules exist but are inert \emph{at the grain at which it is a rock}. Nothing
here claims a grain at which nothing runs.

\begin{proposition}[The observed time is one-dimensional]\label{prop:one-time}
Every observed time is one-dimensional, and one succession is what self-change
lays down: a frame's own changes are strictly ordered and none are concurrent.
Independent successions are not prohibited; the single observed axis is selected,
not forced, and time itself is optional, absent from every inert frame.
\end{proposition}

\begin{remark}[The selection of one time axis is conjectured, not derived]\label{rem:one-time-conjecture}
What selects the single observed time axis is conjectured here to be the same
filter that selects three spatial dimensions --- decoherence and stability ---
and not a necessity. That is a conjecture and is marked as one: unlike the $N=3$
case, which has Ehrenfest's stability result behind it, no comparable result is
cited for the time axis and none is claimed. So one-dimensional time is observed
and selected, not forced; multi-dimensional time is unobserved, not forbidden.
\end{remark}

Space is $N$-dimensional by contingency: the number of positional coordinates the
rules address is rule-content, so ``why $N$'' is ``why this rule-set,'' answered
by self-location (Section~\ref{sec:cosmology}). Why $N=3$ specifically is that
self-location filtered by stability --- Ehrenfest's result that stable orbits and
stable atoms in the usual central-force models are obstructed above three spatial dimensions \citep{ehrenfest1917}, so
observers occur in the three-dimensional rule-sets and we are the observers such a
rule-set produces. And the $3+1$ \emph{signature}, time set apart from space, is
the succession/position distinction: time is the change-axis, the one succession,
space the config-axes, the many positions; one time and $N$ space is
one-succession-and-many-positions read as a signature.

The contingent-not-necessary status of these verdicts --- and of the gauge group,
the dimensionless constants, the arrow, and the value of $c$ --- is licensed by the
minimal-universe argument (Section~\ref{sec:scope}; the mind companion's
Proposition~\ref{MS-prop:minimal-universe}): a rock is a universe at the
grain at which it is a rock, and the minimal universe does without them, so none
can be compulsory.

\section{Identical particles are functional identity}\label{sec:identical}
Graphic Equalisation's Axiom~$-1$ holds that there is no operational view from
nowhere; that two things with no difference that makes a difference are the same
is its \emph{consequence} --- functional identity --- and its
Corollary~\ref{GE-thm:no-frame-free} adds that one cannot say which of two is which
independently of an individuating frame. This is
the quantum treatment of identical particles exactly: two electrons carry no
label, and ``which is which'' names no difference and so has no content.
Classical physics labels identical objects --- treats two identical bodies as
distinct in principle --- and that label is a fiction, a distinction imposed
where none exists.

\begin{proposition}[Quantum statistics, and the Gibbs paradox, are functional identity]\label{prop:gibbs}
Classical Maxwell--Boltzmann counting labels identical particles and over-counts
arrangements; quantum Bose--Einstein and Fermi--Dirac counting treats them as one
where they are indistinguishable. The Gibbs paradox --- the spurious entropy of
mixing two identical gases --- resolves by not labelling them, which is
Axiom~$-1$, and ties to Section~\ref{sec:second-law}: mixing entropy counts
arrangements the frame can tell apart, and an exchange of identical particles is
not one. The $1/N!$ Gibbs correction is functional identity made
quantitative.
\end{proposition}

\begin{remark}[Exclusion needs the fermionic state space]\label{cor:pauli}
Pauli exclusion follows from antisymmetry of fermionic states. Relational
indiscernibility alone does not select it: bosons can occupy the same
single-particle state. The spin--statistics input is required.
\end{remark}

\begin{remark}[The honest limit: the identity argument does not select fermions]\label{rem:spin-stat}
The corollary's argument is weaker than it looks, and the boson case shows why: it applies
verbatim to two bosons in one state, which is observed and which the Gibbs count above needs
to be genuinely two. So Axiom~$-1$ does not \emph{derive} exclusion --- exclusion is
\emph{consistent} with it, and what actually distinguishes the fermion case from the boson
case (the minus sign under exchange, tied to a $2\pi$ rotation) is the spin--statistics
theorem, which needs structure this account does not supply. Nor does ``same in every quantum
number'' by itself individuate: two electrons in distant atoms differ by position, and it is
the full state including that which the argument must range over. The framework grounds
\emph{that} particles are indistinguishable, and hands the fermion/boson split back untouched.
\end{remark}

\section{Positioning and prior art}\label{sec:prior-art}
The pieces are individually established; the contribution is the conjunction
under one change-economy and the mechanisation of shapes prior art postulates,
geometrises, or offers heuristically.

\paragraph{Discrete substrates and emergent spacetime.} Several programmes derive relativistic
structure from a discrete substrate, and two are closer to the change-budget of
Section~\ref{sec:budget} than any continuum account. \citet{darianotosini2013} obtain the metric
from \emph{pure event-counting} on topologically homogeneous causal networks, with time dilation
appearing as an increased density of foliation leaves within one tick of a clock and length
contraction as the corresponding decrease in event density. That is the same move made here and
made earlier, and the counting reading of dilation is theirs, not this paper's. Concretely, where
D'Ariano--Tosini compare the density of foliation leaves associated with a clock tick, this paper
compares how a cycle's variable throughput is divided between motion and internal change: the former
supplies a metric-counting account, while the latter is the additional mechanism claimed for dilation,
mass, and change-load. Three things are
added rather than repeated. Their result is stated for \emph{topologically homogeneous} networks, which
they argue costs no generality since connected events can always be coarse-grained until
homogeneity holds. That defence is sound at the coarse-grained grain, and it is exactly the grain
this paper does not work at: a cycle's length varying with the rules and the graph is the
mechanism here, and coarse-graining until it does not vary would discard the quantity the budget
is about. The difference is therefore real but narrower than it first appears --- not that
homogeneity is a defect in their account, but that the variation it averages over is what carries
the physics in this one. Their counting delivers the metric;
the budget additionally \emph{splits} --- throughput spent on motion or on internal change --- which
is what yields mass as reserved budget, general relativity as change-load rather than only the
special-relativistic metric, and the equivalence principle as the seam between the two
(Section~\ref{sec:sr-gr}). And their result is a result about spacetime; the same economy is run
here through collapse, entanglement and the amplitude weights (Sections~\ref{sec:measurement}
to~\ref{sec:amplitudes}), and belongs to a substrate of which physics is one embodiment rather
than to physics alone.
\citet{knuth2014} derive the Minkowski metric and the Lorentz transformations from consistent
quantification of a causally ordered set with respect to an \emph{embedded observer}, explicitly
without assuming the constancy of $c$; that is the closest existing result to the reading of $c$
given above, and it reaches it by a quantification argument rather than by a throughput economy.
The Wolfram model \citep{wolfram2020,gorard2020rel} derives general covariance from causal
invariance and Lorentz covariance from flat foliations, but stipulates $c$-constancy in its
definition of edge lengths (Section~\ref{sec:c-rate}) and, in the quantum companion
\citep{gorard2020qm}, introduces the Born weights as an a-priori rule rather than recovering them
--- which is the respect in which Section~\ref{sec:amplitudes} differs from it. Causal set theory
\citep{bombelli1987,surya2019} takes order and counting as primitive without rewriting dynamics.
The honest summary is that emergence of relativistic structure from a discrete substrate is
established prior art with several independent routes, and this paper does not claim the emergence.
What it claims is narrower and, so far as these sources go, unoccupied: one budget, split two ways,
carrying the special and general cases together and continuing into the quantum content of the same
substrate without a further postulate at either seam.

\paragraph{Deriving $c$, and mass as internal motion.} That $c$ should be
\emph{derived} rather than postulated has precedent by a different route.
Ignatowsky \citep{ignatowsky1910} and, cleanly, L\'evy-Leblond
\citep{levyleblond1976} obtain the Lorentz transformation from the relativity
principle together with homogeneity, isotropy, and group structure alone, with an
invariant speed \emph{emerging} rather than assumed and its value left empirical;
those derivations use symmetry as the engine where the present one uses the
change-substrate --- $c$ as each frame's own change-throughput unit, $1$ by
construction. The reading of mass as the budget reserved for internal change has
two standing precedents, offered heuristically rather than mechanised:
Schr\"odinger's \emph{zitterbewegung}, the electron's rest energy as internal
motion at $c$ \citep{schrodinger1930}, developed as a rest-mass-from-internal-
circulation interpretation by Hestenes \citep{hestenes1990}; and the elementary
result that confined massless radiation --- light in a box --- carries rest mass
$E/c^{2}$, internalised $c$-motion appearing as inertia. The budget-split itself
is the constant-magnitude four-velocity of Minkowski geometry
\citep{minkowski1908}.

\paragraph{Relativity as postulate.} Einstein's special \citep{einstein1905} and
general \citep{einstein1916} relativity take $c$-invariance and the equivalence
principle as postulates; the present account renders the first a definitional
consequence (Proposition~\ref{prop:c-unit}) and the second a common cause
(Section~\ref{sec:sr-gr}), at the cost of claiming only the shapes.

\paragraph{Observer-first programmes.} The starting point of
Sections~\ref{sec:measurement}--\ref{sec:amplitudes} --- states relative to a frame, no
god's-eye fact --- is shared with several live programmes, and the disagreements matter
more than the agreement.

\citet{muller2020} is the most direct competitor. It takes the first-person observer
state as primitive and adds one postulate --- that algorithmic universal induction fixes
the chances of what the observer sees next --- from which it derives, as an asymptotic
statistical consequence rather than an assumption, that observers find themselves in an
external world governed by simple computable probabilistic laws. The target is exactly
the target here: an objective world \emph{emerging} from an observer-first base rather
than being assumed behind it. The mechanism is incompatible. Solomonoff induction over
sequences of observer states is not graph rewriting under an admissibility grammar, and
nothing in this paper reduces to it, reproduces it, or is supported by it. The two
accounts owe the same debt --- why any shared world at all --- and pay it in different
currency; if M\"uller's route is the right one, the substrate proposed here is at best a
redundant second description of it, and this paper claims no agreement it has not earned.

\citet{hoehn2017} and \citet{hoehnwever2017} are the closest in \emph{shape}. An observer
interrogates a system with binary questions; the system's state \emph{is} the observer's
catalogue of answers, held by the observer and not by the system; and four rules on
information acquisition yield the projective measurement structure, entanglement, and
monogamy. A knowledge structure over a question graph, with constraints governing which
questions may be jointly asked, is typed structure plus admissibility with states
relative to a frame --- nearer to the machinery used here than QBism is, and reaching
qubit quantum theory where Section~\ref{sec:amplitudes} only recovers the Born weight given
contingent rule-content. Worth noting on the other side of the ledger: those four rules yield
the correlation structure of qubits \emph{and rebits} alike, so they do not by themselves select
complex over real. That is independent support for the verdict of Section~\ref{sec:amplitudes}
--- an observer-first reconstruction, pursued to its own conclusion, leaves the same choice
open.

QBism \citep{fuchs2013} shares the observer-first premise and the refusal of
agent-independent states, then goes the opposite way on precisely what this paper cares
about. It declines to derive an objective world, and it treats the Born rule as a
normative addition to Bayesian coherence --- a rule an agent \emph{ought} to obey when
gambling on its own experiences --- rather than as something recovered from deeper
structure. Section~\ref{sec:amplitudes} attempts the recovery, and
Section~\ref{sec:cosmology} keeps the objective world; QBism is therefore the cleanest
available foil, not an ally, and the contrast is the sharpest test of whether the
recovery is worth anything.

Relational quantum mechanics \citep{rovelli1996} is the nearest ancestor of
Section~\ref{sec:measurement}, and the standing problem it has had to answer --- how
different observers' accounts cohere into anything shareable --- is addressed by
cross-perspective links \citep{adlamrovelli2023}, on which one observer's record of
another's outcome is required to agree with that outcome. That is the same problem the
admissibility grammar is asked to solve here and that algorithmic induction solves for
M\"uller. Three different answers to one question; which is right is open.

\paragraph{Constructor theory.} \citet{deutsch2013} proposes re-founding physics on
which physical transformations are \emph{possible} and which \emph{impossible}, and why,
with counterfactuals primitive rather than derived from dynamical laws plus initial
conditions. That is the nearest philosophical relative of the present account: ``which
transformations are admissible, and why'' is close to ``typed graphs plus an
admissibility grammar,'' and both take the possibility structure to be prior to the
trajectory. On the Born rule the relation inverts to competition.
\citet{marletto2016} attempts the same recovery from that admissibility-flavoured base
and arrives somewhere else: probability is not fundamental, and what look like
probabilities emerge from the \emph{unpredictability} of measurement on
superinformation media, giving a non-probabilistic account where
Section~\ref{sec:amplitudes} recovers a probability weight from closure, no-privilege,
and phase-averaging. The two are not quite rival derivations of one rule, since the
explananda differ --- Marletto accounts for the \emph{appearance} of stochasticity in a theory
carrying no probability weight, where Section~\ref{sec:amplitudes} recovers the weight itself ---
and that is a deeper disagreement than a rival derivation would be, not a shallower one. The
honest position is that constructor theory reaches its answer with more formal apparatus and
less contingency than is claimed here, and that this paper's recovery is exposed to it rather
than reinforced by it.

\paragraph{Entropy and the demon.} The observer-relative reading of entropy is
Jaynes's \citep{jaynes1957}; the resolution of Maxwell's demon by the cost of
\emph{erasing} information is Landauer's and Bennett's \citep{landauer1961,bennett1982}; life as
locally paid negative entropy is Schr\"odinger's \citep{schrodinger1944} --- his phrase is \emph{negative entropy}, the contraction
\emph{negentropy} being Brillouin's later coinage. The present
account ties all three to one substrate law --- cost is the write.

\paragraph{Dimensional constants.} That the count of fundamental \emph{dimensional}
constants is conventional is Duff's position in the Duff--Okun--Veneziano exchange
\citep{duff2002trialogue}; the present account gives the reason (they are frame-units) and
draws the boundary at the dimensionless remainder.

\section{Falsifiability commitments}\label{sec:falsifiability}

\noindent\emph{On the tags.} Each commitment below is marked \textsc{[shape]} or \textsc{[detail]}.
A \textsc{[detail]} failing is a revision: the account survives with that mechanism replaced. A
\textsc{[shape]} failing costs \emph{range} rather than a mechanism, because what fails is something
the substrate is claimed to require or to permit --- so the account does not hold where it claimed to,
which is the more serious of the two and is still a boundary rather than an annihilation.
Each is exposed to attack; the scope caveat of Section~\ref{sec:scope} applies
throughout.
\begin{enumerate}[leftmargin=*]
\item \textsc{[detail]}~\textbf{$c$ as frame-unit.} Normalising an accepted unit to one is
conventional. A physical grounding must separately reproduce the observed
invariance of signal speed under its specified transformations. Agreement in
local unit values alone does not test that requirement.
\item \textsc{[detail]}~\textbf{One budget, two dilations.} The budget interpretation must reproduce
proper-time ratios for specified worldlines and metrics, including a clock in a
flat region inside a gravitating shell. Local change-load alone is insufficient;
an explicit rule connecting the budget to the metric is required.
\item \textsc{[shape]}~\textbf{Dimensional/dimensionless boundary.} Fails on exhibition of a
dimensionless constant that is a ratio of nothing --- irreducibly standalone, not expressible as one
quantity against another even in principle --- or if a dimensional constant carries an irreducible
value not absorbable into a frame's units.
\item \textsc{[detail]}~\textbf{The uncertainty principle is the finite-aperture floor.}
A proposed aperture model must construct the conjugate observables and reproduce
their state-preparation uncertainty bound. A bound on measurement disturbance is
a separate claim. Finiteness alone does not establish either bound.
\item \textsc{[detail]}~\textbf{Second-law grounding.} An embodiment must specify entropy,
coarse graining, boundary conditions, and exchange with the environment before
claiming entropy production. Reversible microscopic evolution can accompany an
increase of coarse-grained entropy; overwriting alone is not the second law.
\item \textsc{[detail]}~\textbf{Indistinguishability is functional identity.} A grounding must
recover observable permutation invariance. Fermionic exclusion additionally
requires antisymmetric state space; indistinguishability alone permits bosonic
as well as fermionic statistics.
\item \textsc{[shape]}~\textbf{Measurement is frame-relative collapse.} Fails if a
frame-independent fact about whether collapse has occurred is required to
describe a Wigner's-friend scenario --- if the friend's outcome and the outside
frame's superposition cannot both stand as the event and the map of two frames.
(The Born-rule probabilities are not claimed.)
\item \textsc{[shape]}~\textbf{Entanglement is a superframe dimension.} Fails if the correlation
of an entangled pair is a subframe-local quantity rather than a joint dimension
\textsc{null} in each subframe, or if a reduced single-particle state can be pure
--- i.e.\ if the mixedness is not the inaccessibility of a superframe dimension.
(The Tsirelson bound and correlation strengths are not claimed.)
\item \textsc{[detail]}~\textbf{Conservation is the cost side of symmetry.} Fails if a continuous
\emph{global} symmetry of the action (a residual-free transformation) carries no conserved
quantity --- local gauge redundancies are excluded, carrying constraints rather than charges --- or if a
conserved quantity persists across a genuine change in the rules --- in
particular, if global energy is conserved in a time-dependent (rules-changing)
regime such as cosmic expansion. (The gauge group is claimed to be contingent,
not derived.)
\item \textsc{[detail]}~\textbf{The Past Hypothesis is a seed.} Fails if the low-entropy initial
condition is derived from within --- from the dynamics, with no boundary condition imposed from a
container (Proposition~\ref{prop:past-hyp}). What is claimed is that the seed cannot be derived from
inside, not that its content is beyond reach.
\item \textsc{[detail]}~\textbf{Conditional Born mapping.} The correspondence must construct
states, effects, transformations, and composition satisfying every hypothesis
of its chosen reconstruction theorem. Phase symmetry and linearity alone do
not select the square. The Gleason route separately assumes a normalised,
additive, noncontextual measure in dimension at least three.
\item \textsc{[detail]}~\textbf{Interference is possibility plus superframe forcing.} Fails if
interference requires the particle's own admissibility to select among paths, or
if a which-path coupling that collapses before the screen does not remove the
fringes --- i.e.\ if superposition is not the possibility structure and the
superframe does not apply the combination without passing through the subframe's
grammar.
\item \textsc{[shape]}~\textbf{Time is optional.} Fails if
an inert frame (no firing rules) nonetheless has an in-frame succession.

\item \textsc{[detail]}~\textbf{The one-way speed is an absent quantity.} Fails on a determination of
the one-way speed \emph{in this universe} that rests on no synchrony stipulation. What fails then is
the empirical claim that our rules supply no unlocked counter, not the cycle unit itself: a second
dimension's step count is a different unit and not a finer reading of the same one
(Remark~\ref{rem:oneway-rule}).

\item \textsc{[detail]}~\textbf{The decay tail is required, its exponent is not.} Fails if the
long-time exponent is universal across systems with different threshold structure. The substrate
fixes that there \emph{is} a tail --- exact exponentiality would need spectral support running to
$-\infty$, which the floor forbids --- and cannot fix the exponent, which tracks the density of
states at threshold (Section~\ref{sec:decay}).

\item \textsc{[shape]}~\textbf{In-frame energy is a count of quanta.} Fails if in-frame energy is not
a count of quanta for interacting systems, where the count is not sharp and the physical states are
dressed. Nothing in the spectrum-floor argument survives if that identification fails there, and this
paper does not establish that it holds (Proposition~\ref{prop:spectrum-floor}).

\item \textsc{[detail]}~\textbf{Apparent constancy tracks denominator size.} The constants most
likely to yield a measurable variation are those whose denominators are smallest against the
perturbations available to them, and not the ones deemed most fundamental. Fails if a measured
variation, when found, appears first in constants ordered by depth rather than by that ratio
(Proposition~\ref{prop:dim-boundary}).
\end{enumerate}

\noindent\emph{Self-location} (Proposition~\ref{prop:self-location}) is stated here rather than
counted in the list above, being a consequence rather than a wager: ``why this seed, this constant,
this gauge group'' can be posed only from a trans-frame standpoint that could have obtained a
different one, which Axiom~$-1$ together with frame-relative identity already denies, so what would
have to fail for it to fail is the base axiom and not a risk of its own.

\noindent\emph{Scope note (not a falsifier).} The paper claims shapes only; it would be
overclaiming if any passage were read to derive a field equation, a dimensionless
constant, or a quantum-gravity result --- none is asserted. The Planck-scale claim carries no
separate wager: it is covered as a consequence of the finite-aperture floor together with the
dimensional/dimensionless boundary, and needs no falsifier of its own. That the observed time is $1$-D is
claimed decoherence-selected, not forced --- multi-dimensional time is unobserved, not prohibited;
that space is $N$-D with $N=3$ is likewise contingent and self-located. The minimal-universe
licence (the mind companion's Section~\ref{MS-sec:minimal}), the block universe, free will,
personal identity, and death are treated in that companion and carry their falsifiers there.

\begin{thebibliography}{99}

\bibitem[Adlam and Rovelli(2023)]{adlamrovelli2023}
Emily Adlam and Carlo Rovelli.
\newblock Information is physical: cross-perspective links in relational quantum mechanics.
\newblock \emph{Philosophy of Physics}, 1(1), 2023.

\bibitem[Ali Ahmad et~al.(2022)]{aliahmad2022}
Shadi Ali Ahmad, Thomas D. Galley, Philipp A. H\"ohn, Maximilian P.~E. Lock, and Alexander R.~H. Smith.
\newblock Quantum relativity of subsystems.
\newblock \emph{Physical Review Letters}, 128(17):170401, 2022.

\bibitem[Anderson et~al.(1998)]{anderson1998}
Ronald Anderson, I.~Vetharaniam, and Geoffrey E. Stedman.
\newblock Conventionality of synchronisation, gauge dependence and test theories of relativity.
\newblock \emph{Physics Reports}, 295(3--4):93--180, 1998.

\bibitem[Bell(1964)]{bell1964}
John S. Bell.
\newblock On the Einstein Podolsky Rosen paradox.
\newblock \emph{Physics Physique Fizika}, 1:195--200, 1964.

\bibitem[Bennett(1982)]{bennett1982}
Charles H. Bennett.
\newblock The thermodynamics of computation --- a review.
\newblock \emph{International Journal of Theoretical Physics}, 21(12):905--940, 1982.

\bibitem[BIPM(2019)]{bipm2019}
Bureau International des Poids et Mesures.
\newblock \emph{The International System of Units (SI)}.
\newblock 9th edition, 2019.

\bibitem[Bombelli et~al.(1987)]{bombelli1987}
Luca Bombelli, Joohan Lee, David Meyer, and Rafael D. Sorkin.
\newblock Space-time as a causal set.
\newblock \emph{Physical Review Letters}, 59(5):521--524, 1987.

\bibitem[Bortolotti et~al.(2025)Bortolotti, Curceanu, Di\'osi, Manti, and Piscicchia]{bortolotti2025}
Nicola Bortolotti, Catalina Curceanu, Lajos Di\'osi, Simone Manti, and Kristian Piscicchia.
\newblock Fundamental limits on clock precision from spacetime uncertainty in quantum collapse models.
\newblock \emph{Physical Review Research}, 7:043166, 2025.
\newblock \href{https://doi.org/10.1103/p6tj-lg8l}{doi:10.1103/p6tj-lg8l}.

\bibitem[Brukner(2018)]{brukner2018}
\v{C}aslav Brukner.
\newblock A no-go theorem for observer-independent facts.
\newblock \emph{Entropy}, 20(5):350, 2018.

\bibitem[Chiribella et~al.(2011)]{chiribella2011}
Giulio Chiribella, Giacomo Mauro D'Ariano, and Paolo Perinotti.
\newblock Informational derivation of quantum theory.
\newblock \emph{Physical Review A}, 84(1):012311, 2011.

\bibitem[D'Ariano and Tosini(2013)]{darianotosini2013}
Giacomo Mauro D'Ariano and Alessandro Tosini.
\newblock Emergence of space--time from topologically homogeneous causal networks.
\newblock \emph{Studies in History and Philosophy of Modern Physics}, 44(3):294--299, 2013.

\bibitem[Deutsch(2013)]{deutsch2013}
David Deutsch.
\newblock Constructor theory.
\newblock \emph{Synthese}, 190(18):4331--4359, 2013.

\bibitem[Duff(2002)]{duff2002}
Michael J. Duff.
\newblock Comment on time-variation of fundamental constants.
\newblock arXiv:hep-th/0208093, 2002.

\bibitem[Duff et~al.(2002)]{duff2002trialogue}
Michael J. Duff, Lev B. Okun, and Gabriele Veneziano.
\newblock Trialogue on the number of fundamental constants.
\newblock \emph{Journal of High Energy Physics}, 03:023, 2002. arXiv:physics/0110060.

\bibitem[Ehrenfest(1917)]{ehrenfest1917}
Paul Ehrenfest.
\newblock In what way does it become manifest in the fundamental laws of physics that space has three dimensions?
\newblock \emph{Proceedings of the Amsterdam Academy}, 20:200--209, 1917.

\bibitem[Einstein(1905)]{einstein1905}
Albert Einstein.
\newblock Zur Elektrodynamik bewegter K\"orper.
\newblock \emph{Annalen der Physik}, 322(10):891--921, 1905.

\bibitem[Einstein(1916)]{einstein1916}
Albert Einstein.
\newblock Die Grundlage der allgemeinen Relativit\"atstheorie.
\newblock \emph{Annalen der Physik}, 354(7):769--822, 1916.

\bibitem[Frauchiger and Renner(2018)]{frauchiger2018}
Daniela Frauchiger and Renato Renner.
\newblock Quantum theory cannot consistently describe the use of itself.
\newblock \emph{Nature Communications}, 9:3711, 2018.

\bibitem[Fuchs and Schack(2013)]{fuchs2013}
Christopher A. Fuchs and R\"udiger Schack.
\newblock Quantum-Bayesian coherence.
\newblock \emph{Reviews of Modern Physics}, 85(4):1693--1715, 2013.

\bibitem[Gleason(1957)]{gleason1957}
Andrew M. Gleason.
\newblock Measures on the closed subspaces of a Hilbert space.
\newblock \emph{Journal of Mathematics and Mechanics}, 6:885--893, 1957.

\bibitem[Gorard(2020a)]{gorard2020rel}
Jonathan Gorard.
\newblock Some relativistic and gravitational properties of the Wolfram model.
\newblock \emph{Complex Systems}, 29(2):599--654, 2020.

\bibitem[Gorard(2020b)]{gorard2020qm}
Jonathan Gorard.
\newblock Some quantum mechanical properties of the Wolfram model.
\newblock \emph{Complex Systems}, 29(2):537--598, 2020.

\bibitem[Hardy(2001)]{hardy2001}
Lucien Hardy.
\newblock Quantum theory from five reasonable axioms.
\newblock arXiv:quant-ph/0101012, 2001.

\bibitem[Hau et~al.(1999)]{hau1999}
Lene Vestergaard Hau, S.~E. Harris, Zachary Dutton, and Cyrus H. Behroozi.
\newblock Light speed reduction to 17 metres per second in an ultracold atomic gas.
\newblock \emph{Nature}, 397(6720):594--598, 1999.

\bibitem[Hensen et~al.(2015)]{hensen2015}
B.~Hensen, H.~Bernien, A.~E. Dr\'eau, A.~Reiserer, N.~Kalb, M.~S. Blok, J.~Ruitenberg, R.~F.~L. Vermeulen, R.~N. Schouten, C.~Abell\'an, W.~Amaya, V.~Pruneri, M.~W. Mitchell, M.~Markham, D.~J. Twitchen, D.~Elkouss, S.~Wehner, T.~H. Taminiau, and R.~Hanson.
\newblock Loophole-free {B}ell inequality violation using electron spins separated by 1.3 kilometres.
\newblock \emph{Nature}, 526(7575):682--686, 2015.

\bibitem[Hestenes(1990)]{hestenes1990}
David Hestenes.
\newblock The zitterbewegung interpretation of quantum mechanics.
\newblock \emph{Foundations of Physics}, 20(10):1213--1232, 1990.

\bibitem[H\"ohn(2017)]{hoehn2017}
Philipp A. H\"ohn.
\newblock Toolbox for reconstructing quantum theory from rules on information acquisition.
\newblock \emph{Quantum}, 1:38, 2017.

\bibitem[H\"ohn and Wever(2017)]{hoehnwever2017}
Philipp A. H\"ohn and Christopher S.~P. Wever.
\newblock Quantum theory from questions.
\newblock \emph{Physical Review A}, 95(1):012102, 2017.

\bibitem[Ignatowsky(1910)]{ignatowsky1910}
Wladimir~A. von Ignatowsky.
\newblock Einige allgemeine Bemerkungen zum Relativit\"atsprinzip.
\newblock \emph{Physikalische Zeitschrift}, 11:972--976, 1910.

\bibitem[Jaynes(1957)]{jaynes1957}
Edwin T. Jaynes.
\newblock Information theory and statistical mechanics.
\newblock \emph{Physical Review}, 106:620--630, 1957.

\bibitem[Khalfin(1957)]{khalfin1957}
Leonid A. Khalfin.
\newblock Contribution to the decay theory of a quasi-stationary state.
\newblock Russian original 1957; \emph{Soviet Physics JETP}, 6(6):1053--1063, 1958.

\bibitem[Knuth and Bahreyni(2014)]{knuth2014}
Kevin H. Knuth and Newshaw Bahreyni.
\newblock A potential foundation for emergent space-time.
\newblock \emph{Journal of Mathematical Physics}, 55(11):112501, 2014.

\bibitem[Kochen and Specker(1967)]{kochen1967}
Simon Kochen and Ernst P. Specker.
\newblock The problem of hidden variables in quantum mechanics.
\newblock \emph{Journal of Mathematics and Mechanics}, 17:59--87, 1967.

\bibitem[Landauer(1961)]{landauer1961}
Rolf Landauer.
\newblock Irreversibility and heat generation in the computing process.
\newblock \emph{IBM Journal of Research and Development}, 5(3):183--191, 1961.

\bibitem[L\'evy-Leblond(1976)]{levyleblond1976}
Jean-Marc L\'evy-Leblond.
\newblock One more derivation of the Lorentz transformation.
\newblock \emph{American Journal of Physics}, 44(3):271--277, 1976.

\bibitem[Malament(1977)]{malament1977}
David B. Malament.
\newblock Causal theories of time and the conventionality of simultaneity.
\newblock \emph{No\^us}, 11(3):293--300, 1977.

\bibitem[Marletto(2016)]{marletto2016}
Chiara Marletto.
\newblock Constructor theory of probability.
\newblock \emph{Proceedings of the Royal Society A}, 472(2192):20150883, 2016.

\bibitem[Masanes et~al.(2019)Masanes, Galley, and M\"uller]{masanes2019}
Llu\'is Masanes, Thomas D. Galley, and Markus P. M\"uller.
\newblock The measurement postulates of quantum mechanics are operationally redundant.
\newblock \emph{Nature Communications}, 10:1361, 2019.
\newblock \url{https://doi.org/10.1038/s41467-019-09348-x}.

\bibitem[Masanes and M\"uller(2011)]{masanes2011}
Llu\'is Masanes and Markus P. M\"uller.
\newblock A derivation of quantum theory from physical requirements.
\newblock \emph{New Journal of Physics}, 13(6):063001, 2011.

\bibitem[Minkowski(1908)]{minkowski1908}
Hermann Minkowski.
\newblock Raum und Zeit.
\newblock Address delivered 1908; \emph{Physikalische Zeitschrift}, 10:104--111, 1909.

\bibitem[M\"uller(2020)]{muller2020}
Markus P. M\"uller.
\newblock Law without law: from observer states to physics via algorithmic information theory.
\newblock \emph{Quantum}, 4:301, 2020.

\bibitem[M\"uller(2021)]{muller2021}
Markus P. M\"uller.
\newblock Probabilistic theories and reconstructions of quantum theory.
\newblock \emph{SciPost Physics Lecture Notes}, 28, 2021.

\bibitem[Reichenbach(1958)]{reichenbach1958}
Hans Reichenbach.
\newblock \emph{The Philosophy of Space and Time}.
\newblock Dover, New York, 1958.
\newblock Translation of \emph{Philosophie der Raum-Zeit-Lehre}, 1928.

\bibitem[Rissanen(1978)]{rissanen1978}
Jorma Rissanen.
\newblock Modeling by shortest data description.
\newblock \emph{Automatica}, 14(5):465--471, 1978.

\bibitem[Rothe et~al.(2006)]{rothe2006}
Carsten Rothe, Sergey I. Hintschich, and Andrew P. Monkman.
\newblock Violation of the exponential-decay law at long times.
\newblock \emph{Physical Review Letters}, 96(16):163601, 2006.

\bibitem[Rovelli(1996)]{rovelli1996}
Carlo Rovelli.
\newblock Relational quantum mechanics.
\newblock \emph{International Journal of Theoretical Physics}, 35:1637--1678, 1996.

\bibitem[Schlosshauer and Fine(2005)]{schlosshauerfine2005}
Maximilian Schlosshauer and Arthur Fine.
\newblock On Zurek's derivation of the Born rule.
\newblock \emph{Foundations of Physics}, 35:197--213, 2005.
\newblock \url{https://arxiv.org/abs/quant-ph/0312058}.

\bibitem[Schr\"odinger(1930)]{schrodinger1930}
Erwin Schr\"odinger.
\newblock \"Uber die kr\"aftefreie Bewegung in der relativistischen Quantenmechanik.
\newblock \emph{Sitzungsberichte der Preussischen Akademie der Wissenschaften, Physikalisch-mathematische Klasse}, 24:418--428, 1930.

\bibitem[Schr\"odinger(1944)]{schrodinger1944}
Erwin Schr\"odinger.
\newblock \emph{What Is Life? The Physical Aspect of the Living Cell}.
\newblock Cambridge University Press, Cambridge, 1944.

\bibitem[Selby et~al.(2021)]{selby2021}
John H. Selby, Carlo Maria Scandolo, and Bob Coecke.
\newblock Reconstructing quantum theory from diagrammatic postulates.
\newblock \emph{Quantum}, 5:445, 2021.

\bibitem[Streater and Wightman(1964)]{streater1964}
Raymond F. Streater and Arthur S. Wightman.
\newblock \emph{PCT, Spin and Statistics, and All That}.
\newblock W. A. Benjamin, New York, 1964.

\bibitem[Surya(2019)]{surya2019}
Sumati Surya.
\newblock The causal set approach to quantum gravity.
\newblock \emph{Living Reviews in Relativity}, 22:5, 2019.

\bibitem[Uzan(2011)]{uzan2011}
Jean-Philippe Uzan.
\newblock Varying constants, gravitation and cosmology.
\newblock \emph{Living Reviews in Relativity}, 14:2, 2011.

\bibitem[Webb et~al.(2011)]{webb2011}
John K. Webb, Julian A. King, Michael T. Murphy, Victor V. Flambaum, Robert F. Carswell, and
Matthew B. Bainbridge.
\newblock Indications of a spatial variation of the fine structure constant.
\newblock \emph{Physical Review Letters}, 107(19):191101, 2011.

\bibitem[Wilkinson et~al.(1997)]{wilkinson1997}
Steven R. Wilkinson, Cyrus F. Bharucha, Martin C. Fischer, Kirk W. Madison, Patrick R. Morrow,
Qian Niu, Bala Sundaram, and Mark G. Raizen.
\newblock Experimental evidence for non-exponential decay in quantum tunnelling.
\newblock \emph{Nature}, 387(6633):575--577, 1997.

\bibitem[Wolfram(2020)]{wolfram2020}
Stephen Wolfram.
\newblock A class of models with the potential to represent fundamental physics.
\newblock \emph{Complex Systems}, 29(2):107--536, 2020.

\bibitem[Zurek(2003)]{zurek2003}
Wojciech H. Zurek.
\newblock Decoherence, einselection, and the quantum origins of the classical.
\newblock \emph{Reviews of Modern Physics}, 75(3):715--775, 2003.

\end{thebibliography}
\end{document}
